\documentclass[12pt,a4paper,openright]{report}

% --- Encodage et langue ---
\usepackage[utf8]{inputenc}
\usepackage[T1]{fontenc}
\usepackage[french]{babel}

% --- Mathématiques ---
\usepackage{amsmath, amssymb, amsthm, amsfonts}

% --- Mise en page ---
\usepackage{geometry}
\geometry{hmargin=2.5cm,vmargin=2.5cm}
\usepackage{xcolor}
\usepackage{tikz}

% --- Code source (Python) ---
\usepackage{listings}
\lstset{
    language=Python,
    basicstyle=\ttfamily\small,
    keywordstyle=\color{blue}\bfseries,
    stringstyle=\color{red},
    commentstyle=\color{green!60!black}\itshape,
    numbers=left,
    numberstyle=\tiny\color{gray},
    stepnumber=1,
    frame=single,
    breaklines=true
}

% --- Liens cliquables (table des matières navigable) ---
\usepackage[colorlinks=true,
            linkcolor=blue!60!black,
            urlcolor=blue!60!black,
            citecolor=blue!60!black,
            bookmarks=true,
            bookmarksnumbered=true,
            pdftitle={Algèbre générale - ECUE111},
            pdfauthor={Aymen Ben Brik}]{hyperref}

% --- Environnements de théorèmes (numérotation par section, continue) ---
\newtheorem{theoreme}{Théorème}[section]
\newtheorem{proposition}[theoreme]{Proposition}
\newtheorem{lemme}[theoreme]{Lemme}
\newtheorem{corollaire}[theoreme]{Corollaire}
\theoremstyle{definition}
\newtheorem{definition}{Définition}[section]
\newtheorem{exemple}{Exemple}[section]
\theoremstyle{remark}
\newtheorem{remarque}{Remarque}[section]
\newtheorem{propriete}{Propriété}[section]

% --- Métadonnées ---
\title{
    \vspace*{2cm}
    {\Huge \textbf{Algèbre générale}} \\[0.5cm]
    {\Large Module ECUE111 --- S1, 2025--2026} \\[1cm]
    {\large Filière Mathématiques --- Option Sciences des données} \\[0.3cm]
    {\large 1\textsuperscript{ère} année}
}
\author{Aymen Ben Brik}
\date{\today}

\begin{document}

% --- Page de titre globale ---
\begin{titlepage}
\maketitle
\vfill
\begin{center}
\begin{tabular}{ll}
\textbf{Code UE} & UEF110 \\
\textbf{Code ECUE} & ECUE111 \\
\textbf{Volume horaire} & 84 h présentiel (42 h cours + 42 h TD) \\
\textbf{Coefficient} & 3.5 \\
\textbf{Crédits ECTS} & 7 \\
\textbf{Régime} & CC mixte \\
\textbf{Responsable du module} & Moufid Chaima \\
\textbf{Rédaction du support} & Aymen Ben Brik \\
\end{tabular}
\end{center}
\vfill
\end{titlepage}

% --- Avant-propos ---
\chapter*{Avant-propos}
\addcontentsline{toc}{chapter}{Avant-propos}

Ce document constitue le support complet du module \emph{Algèbre générale} (ECUE111), en première année de la filière Mathématiques (option Sciences des données). Il couvre les six chapitres du programme officiel :
\begin{enumerate}
    \item Calculs algébriques (sommes, produits, formule du binôme)
    \item Vocabulaire ensembliste (logique, ensembles, dénombrement, relations)
    \item Arithmétique dans $\mathbb{Z}$
    \item Structures algébriques (groupes, anneaux, corps)
    \item Polynômes
    \item Fractions rationnelles
\end{enumerate}

\paragraph{Pédagogie.} Chaque chapitre suit une structure unifiée : une \emph{motivation} concrète (en informatique, science des données, cryptographie ou ingénierie), une \emph{approche par problème}, des \emph{définitions et théorèmes} avec démonstrations, des \emph{exemples résolus} à difficulté ascendante, du \emph{code Python} (\texttt{sympy}, \texttt{numpy}, \texttt{itertools}) et des \emph{exercices à faire}. Chaque chapitre se conclut par un \emph{TP Python structuré} (annotations Bloom + difficulté), construisant progressivement des bibliothèques réutilisables.

\paragraph{Acquis d'apprentissage.}
\begin{itemize}
    \item \textbf{AA1} : Maîtriser les opérations algébriques fondamentales et leurs propriétés.
    \item \textbf{AA2} : Utiliser logique et théorie des ensembles pour structurer un raisonnement.
    \item \textbf{AA3} : Appliquer les principes d'arithmétique (division euclidienne, premiers, congruences).
    \item \textbf{AA4} : Manipuler les polynômes (opérations, factorisation, racines).
    \item \textbf{AA5} : Manipuler les fractions rationnelles et simplifications algébriques.
    \item \textbf{AA6} : Rédiger et justifier des arguments mathématiques de manière rigoureuse et claire.
\end{itemize}

\paragraph{Code source.} Le repo GitHub \texttt{Aymenbenbrik/Algebre1-2026} contient les fichiers \LaTeX{} sources de chaque chapitre, ainsi que les notes SmartBoard des séances.

% --- Table des matières ---
\tableofcontents

% --- Chapitres ---
\input{ch1}
\input{ch2}
% On ajuste le compteur pour que ce soit le Chapitre 3

\chapter{Rappels d'arithmétique dans l’ensemble des entiers relatifs}

% Introduction brève du chapitre (optionnel)
Ce chapitre a pour objectif de consolider les bases de l'arithmétique dans l'ensemble des entiers relatifs ($\mathbb{Z}$).

% Inclusion des différentes sections
\section{Division euclidienne, congruence}

\subsection{Motivation}
En informatique et en sciences de l'ingénieur, les nombres entiers ne sont pas de simples compteurs. L'arithmétique modulaire, qui repose sur la division euclidienne et les congruences, est le moteur de la cryptographie moderne (comme les protocoles qui sécurisent nos communications), mais aussi des algorithmes de hachage et de la génération de nombres pseudo-aléatoires. Maîtriser ces concepts permet de comprendre comment les machines traitent, transforment et sécurisent l'information de manière cyclique.

\subsection{Approche par problème}
Imaginons un capteur dans un système industriel qui enregistre une mesure de température toutes les 7 heures. Si le capteur est activé un lundi à 00h00, quel jour de la semaine la 100\textsuperscript{ème} mesure sera-t-elle prise ? 

Ce type de problème cyclique ne nécessite pas de calculer le temps total écoulé de manière absolue, mais plutôt de s'intéresser au \textit{reste} de la division du temps total par la durée d'un cycle. C'est précisément l'objet de la division euclidienne et de la notion de congruence, qui permettent de simplifier les grands nombres en se concentrant uniquement sur leur comportement périodique.

\subsection{Définitions et théorèmes}

\subsubsection{La division euclidienne}

\begin{theoreme}[Division euclidienne dans $\mathbb{Z}$]
Soit $a \in \mathbb{Z}$ et $b \in \mathbb{Z}^*$. Il existe un unique couple d'entiers $(q, r) \in \mathbb{Z} \times \mathbb{N}$ tel que :
$$a = bq + r \quad \text{avec} \quad 0 \leq r < |b|$$
\end{theoreme}

\begin{remarque}
Dans cette écriture, $a$ est appelé le \textbf{dividende}, $b$ le \textbf{diviseur}, $q$ le \textbf{quotient} et $r$ le \textbf{reste}. Il est crucial de noter que le reste $r$ est toujours \textbf{positif ou nul}, même si $a$ ou $b$ sont négatifs.
\end{remarque}

\subsubsection{Congruences}

La notion de congruence permet de formaliser mathématiquement ces phénomènes cycliques en travaillant directement sur les restes.

\begin{definition}[Congruence modulo $n$]
Soit $n \in \mathbb{N}^*$. Deux entiers relatifs $a$ et $b$ sont dits congrus modulo $n$ si et seulement si $n$ divise leur différence $a - b$. On note :
$$a \equiv b \pmod{n}$$
Ceci est équivalent à dire que $a$ et $b$ ont le même reste dans la division euclidienne par $n$.
\end{definition}

\begin{propriete}[Opérations sur les congruences]
La relation de congruence est une relation d'équivalence compatible avec l'addition et la multiplication. Si $a \equiv b \pmod{n}$ et $c \equiv d \pmod{n}$, alors :
\begin{itemize}
    \item $a + c \equiv b + d \pmod{n}$
    \item $ac \equiv bd \pmod{n}$
    \item Pour tout entier naturel $k$, $a^k \equiv b^k \pmod{n}$
\end{itemize}
\end{propriete}

\subsection{Exemples de difficulté croissante}

\begin{exemple}[Division avec des entiers négatifs]
Effectuer la division euclidienne de $a = -17$ par $b = 5$.
\end{exemple}

\textbf{Correction :} On cherche $q \in \mathbb{Z}$ et $r \in \mathbb{N}$ tels que $-17 = 5q + r$ avec $0 \leq r < 5$. 
Si l'on prend $q = -3$, on a $5 \times (-3) = -15$. Le reste serait de $-2$, ce qui est impossible car le reste doit être positif. 
On doit donc « descendre » au multiple inférieur : on prend $q = -4$. 
On obtient : $-17 = 5 \times (-4) + 3$. 
Le quotient est donc $q = -4$ et le reste est $r = 3$.

\begin{exemple}[Calcul modulaire sur de grandes puissances]
Déterminer le reste de la division euclidienne de $3^{2026}$ par $7$.
\end{exemple}

\textbf{Correction :} L'objectif est de trouver une puissance de $3$ qui soit congrue à $1$ ou $-1$ modulo $7$.
Calculons les premières puissances :
\begin{itemize}
    \item $3^1 \equiv 3 \pmod{7}$
    \item $3^2 = 9 \equiv 2 \pmod{7}$
    \item $3^3 = 27 \equiv -1 \pmod{7}$
\end{itemize}
C'est ce qu'il nous faut ! Exprimons l'exposant $2026$ en fonction de $3$. La division euclidienne de $2026$ par $3$ donne : $2026 = 3 \times 675 + 1$.
On peut donc écrire :
$$3^{2026} = 3^{3 \times 675 + 1} = (3^3)^{675} \times 3^1$$
En passant aux congruences modulo $7$ :
$$3^{2026} \equiv (-1)^{675} \times 3 \pmod{7}$$
Comme $675$ est impair, $(-1)^{675} = -1$.
$$3^{2026} \equiv -1 \times 3 \pmod{7} \equiv -3 \pmod{7}$$
Puisque le reste doit être positif, on ajoute $7$ : $-3 + 7 = 4$.
$$3^{2026} \equiv 4 \pmod{7}$$
Le reste de la division euclidienne de $3^{2026}$ par $7$ est donc $4$.

\subsection{Exercice pratique avec Python}

\textbf{Objectif :} Implémenter la division euclidienne et vérifier des congruences.

En Python, les opérateurs \texttt{//} (quotient) et \texttt{\%} (modulo) calculent la division euclidienne en respectant la règle mathématique du reste positif, même pour les entiers négatifs. C'est un avantage majeur par rapport à d'autres langages comme le C ou le Java.

\begin{verbatim}
def division_euclidienne(a, b):
    """Calcule et affiche le quotient et le reste de a par b."""
    if b == 0:
        return "Erreur : division par zéro"
    q = a // b
    r = a % b
    print(f"Dividende {a} = {b} * {q} + {r}")
    return q, r

def sont_congrus(a, b, n):
    """Vérifie si a et b sont congrus modulo n."""
    return (a % n) == (b % n)

# --- Tests ---
print("--- Division euclidienne ---")
division_euclidienne(17, 5)
division_euclidienne(-17, 5) # Python gère le reste positif !

print("\n--- Congruences ---")
# Vérifions le résultat de l'exemple 2 : 3^2026 = 4 mod 7
# En Python, 3**2026 calcule la puissance exacte instantanément.
resultat = sont_congrus(3**2026, 4, 7)
print(f"3^2026 est-il congru à 4 modulo 7 ? {resultat}")
\end{verbatim}


\section{PGCD et PPCM}

\subsection{Motivation}
Dans de nombreux domaines de l'ingénierie et de l'informatique, nous sommes confrontés à des problèmes d'optimisation et de synchronisation. Comment diviser une surface de stockage en blocs carrés les plus grands possibles sans perte d'espace ? Comment déterminer le moment où deux processus exécutés en parallèle, avec des temps de cycle différents, vont se superposer ? Le PGCD (Plus Grand Commun Diviseur) et le PPCM (Plus Petit Commun Multiple) sont les outils mathématiques fondamentaux pour répondre à ces questions.

\subsection{Approche par problème}
Imaginons deux tâches dans un système d'exploitation embarqué. La tâche A s'exécute toutes les 15 millisecondes et la tâche B toutes les 25 millisecondes. Si elles démarrent exactement en même temps à $t=0$, au bout de combien de temps s'exécuteront-elles à nouveau simultanément ?

Pour résoudre cela, on cherche un multiple commun à 15 et 25, et idéalement le plus petit possible pour anticiper le premier conflit : c'est le PPCM. 
À l'inverse, si l'on souhaite diviser une mémoire de 1500 octets et une autre de 2500 octets en segments de taille identique, la plus grande taille possible pour ces segments sera donnée par le PGCD.

\subsection{Définitions et théorèmes}

\subsubsection{Le Plus Grand Commun Diviseur (PGCD)}

\begin{definition}[PGCD]
Soient $a$ et $b$ deux entiers relatifs non tous les deux nuls. L'ensemble des diviseurs communs à $a$ et $b$ est une partie finie et non vide de $\mathbb{Z}$. Le plus grand élément de cet ensemble est un entier naturel appelé Plus Grand Commun Diviseur de $a$ et $b$. 
On le note $\text{PGCD}(a, b)$ ou simplement $a \wedge b$.
\end{definition}

\begin{definition}[Nombres premiers entre eux]
Deux entiers relatifs $a$ et $b$ sont dits premiers entre eux si et seulement si leur seul diviseur commun positif est 1, c'est-à-dire :
$$\text{PGCD}(a, b) = 1$$
\end{definition}

\subsubsection{Le Plus Petit Commun Multiple (PPCM)}

\begin{definition}[PPCM]
Soient $a$ et $b$ deux entiers relatifs non nuls. L'ensemble des multiples communs strictement positifs de $a$ et $b$ possède un plus petit élément. Cet entier naturel est appelé Plus Petit Commun Multiple de $a$ et $b$.
On le note $\text{PPCM}(a, b)$ ou simplement $a \vee b$.
\end{definition}

\subsubsection{Relation fondamentale}

Il existe un lien analytique très fort entre le PGCD et le PPCM de deux nombres.

\begin{theoreme}[Lien entre PGCD et PPCM]
Pour tous entiers relatifs non nuls $a$ et $b$, le produit de leur PGCD et de leur PPCM est égal à la valeur absolue de leur produit :
$$\text{PGCD}(a, b) \times \text{PPCM}(a, b) = |ab|$$
\end{theoreme}

\subsection{Exemples de difficulté croissante}

\begin{exemple}[Calcul par décomposition en facteurs premiers]
Déterminer le PGCD et le PPCM de $a = 60$ et $b = 84$.
\end{exemple}

\textbf{Correction :} On décompose d'abord en produit de facteurs premiers.
$60 = 2^2 \times 3^1 \times 5^1$
$84 = 2^2 \times 3^1 \times 7^1$

\begin{itemize}
    \item Pour le \textbf{PGCD}, on prend les facteurs premiers \textit{communs} avec la \textit{plus petite} puissance :
    $$\text{PGCD}(60, 84) = 2^2 \times 3^1 = 12$$
    \item Pour le \textbf{PPCM}, on prend \textit{tous} les facteurs premiers présents avec la \textit{plus grande} puissance :
    $$\text{PPCM}(60, 84) = 2^2 \times 3^1 \times 5^1 \times 7^1 = 420$$
\end{itemize}
On peut vérifier le théorème : $12 \times 420 = 5040$ et $60 \times 84 = 5040$.

\begin{exemple}[Démonstration algébrique]
Montrer que pour tout entier naturel $n$, les nombres $a = n$ et $b = n+1$ sont premiers entre eux.
\end{exemple}

\textbf{Correction :} Soit $d$ un diviseur commun positif à $n$ et $n+1$. 
Par définition, il existe des entiers $k_1$ et $k_2$ tels que $n = d \cdot k_1$ et $n+1 = d \cdot k_2$.
Si $d$ divise $n$ et $n+1$, il divise aussi toute combinaison linéaire de ces deux nombres, et en particulier leur différence :
$$(n+1) - n = d \cdot k_2 - d \cdot k_1 = d(k_2 - k_1)$$
Donc $1 = d(k_2 - k_1)$. Le seul entier naturel qui divise $1$ est $1$. Donc $d = 1$.
Leur seul diviseur commun étant 1, on a bien $\text{PGCD}(n, n+1) = 1$. Ils sont premiers entre eux.

\subsection{Exercice pratique avec Python}

\textbf{Objectif :} Utiliser la bibliothèque \texttt{math} de Python pour calculer le PGCD et déduire le PPCM.

Depuis Python 3.9, les fonctions \texttt{gcd} (Greatest Common Divisor) et \texttt{lcm} (Least Common Multiple) sont directement incluses dans le module mathématique standard.

\begin{verbatim}
import math

def calculer_pgcd_ppcm(a, b):
    """Affiche le PGCD et le PPCM de deux entiers."""
    if a == 0 and b == 0:
        return "Le PGCD de 0 et 0 n'est pas défini."
        
    # Utilisation des fonctions natives de Python
    pgcd_val = math.gcd(a, b)
    
    # Si on utilise une version de Python < 3.9, on peut déduire le PPCM :
    # ppcm_val = abs(a * b) // pgcd_val
    # Sinon, on utilise math.lcm :
    ppcm_val = math.lcm(a, b)
    
    print(f"Pour a={a} et b={b} :")
    print(f" -> PGCD = {pgcd_val}")
    print(f" -> PPCM = {ppcm_val}")
    print(f" -> Vérification PGCD * PPCM = {pgcd_val * ppcm_val} (|ab| = {abs(a*b)})")

# --- Tests ---
calculer_pgcd_ppcm(60, 84)
print("-" * 20)
calculer_pgcd_ppcm(15, 25) # Retour à notre problème de synchronisation
\end{verbatim}
\section{Théorème de Gauss, Identité de Bézout, Algorithme d’Euclide}

\subsection{Motivation}
La théorie des nombres ne se contente pas d'étudier les propriétés abstraites des entiers ; elle fournit les algorithmes qui sécurisent le monde numérique. L'algorithme d'Euclide est l'un des plus vieux algorithmes au monde, et sa version "étendue" (liée à l'identité de Bézout) permet de calculer des inverses modulaires. Ces calculs sont le cœur mathématique du chiffrement RSA, utilisé quotidiennement pour sécuriser les transactions bancaires et les communications sur Internet.

\subsection{Approche par problème}
Imaginez que vous disposiez d'un récipient de 5 litres et d'un autre de 3 litres, sans aucune graduation. Comment faire pour obtenir exactement 1 litre d'eau ? 

Mathématiquement, cela revient à chercher deux entiers $u$ et $v$ (le nombre de fois où l'on remplit ou vide chaque récipient) tels que $5u + 3v = 1$. Cette équation, appelée équation diophantienne linéaire, possède des solutions entières uniquement grâce aux relations profondes qui lient le nombre 5, le nombre 3, et leur PGCD. Le théorème de Bézout et l'algorithme d'Euclide nous donnent la méthode systématique pour trouver ces solutions.

\subsection{Définitions et théorèmes}

\subsubsection{L'Algorithme d'Euclide}

L'algorithme d'Euclide permet de calculer le PGCD de deux entiers sans avoir à chercher tous leurs diviseurs. Il repose sur le lemme suivant : si $a = bq + r$, alors $\text{PGCD}(a, b) = \text{PGCD}(b, r)$.

\begin{theoreme}[Algorithme d'Euclide]
Soient $a$ et $b$ deux entiers naturels non nuls avec $a > b$. En effectuant des divisions euclidiennes successives :
\begin{align*}
a &= bq_1 + r_1 \quad (0 < r_1 < b) \\
b &= r_1q_2 + r_2 \quad (0 < r_2 < r_1) \\
r_1 &= r_2q_3 + r_3 \quad (0 < r_3 < r_2) \\
&\dots \\
r_{n-2} &= r_{n-1}q_n + r_n \quad (0 < r_n < r_{n-1}) \\
r_{n-1} &= r_nq_{n+1} + 0
\end{align*}
Le dernier reste non nul, $r_n$, est exactement le $\text{PGCD}(a, b)$.
\end{theoreme}



[Image of flowchart of the Euclidean algorithm]


\subsubsection{Identité de Bézout}

\begin{theoreme}[Identité de Bézout]
Deux entiers relatifs $a$ et $b$ sont premiers entre eux si et seulement s'il existe deux entiers relatifs $u$ et $v$ tels que :
$$au + bv = 1$$
\end{theoreme}

\begin{propriete}[Généralisation - Théorème de Bachet-Bézout]
Soient $a$ et $b$ deux entiers non nuls. Si $d = \text{PGCD}(a, b)$, alors il existe deux entiers relatifs $u$ et $v$ tels que :
$$au + bv = d$$
De plus, l'équation $ax + by = c$ admet des solutions entières $(x, y)$ si et seulement si $c$ est un multiple de $d$.
\end{propriete}

\subsubsection{Théorème de Gauss}

\begin{theoreme}[Théorème de Gauss]
Soient $a$, $b$ et $c$ trois entiers relatifs non nuls. 
Si $a$ divise le produit $bc$ et si $a$ et $b$ sont premiers entre eux ($\text{PGCD}(a, b) = 1$), alors $a$ divise $c$.
\end{theoreme}

\subsection{Exemples de difficulté croissante}

\begin{exemple}[Recherche de PGCD par l'algorithme d'Euclide]
Déterminer le PGCD de $420$ et $135$.
\end{exemple}

\textbf{Correction :} On applique l'algorithme d'Euclide par divisions successives :
\begin{itemize}
    \item $420 = 135 \times 3 + 15$ (le reste est 15)
    \item $135 = 15 \times 9 + 0$ (le reste est 0)
\end{itemize}
Le dernier reste non nul est $15$. Donc, $\text{PGCD}(420, 135) = 15$.

\begin{exemple}[Résolution d'une équation avec Bézout et Gauss]
Résoudre dans $\mathbb{Z}^2$ l'équation $(E) : 5x - 3y = 2$.
\end{exemple}

\textbf{Correction :} 
1. \textbf{Recherche d'une solution particulière :} On remarque de manière évidente que $5 \times 1 - 3 \times 1 = 2$. Donc le couple $(1, 1)$ est une solution particulière.
2. \textbf{Soustraction :} Soit $(x, y)$ une solution. On a :
$$5x - 3y = 2$$
$$5(1) - 3(1) = 2$$
En soustrayant ligne à ligne : $5(x - 1) - 3(y - 1) = 0$, soit $5(x - 1) = 3(y - 1)$.
3. \textbf{Application du théorème de Gauss :} $5$ divise $3(y - 1)$. Or $\text{PGCD}(5, 3) = 1$. D'après le théorème de Gauss, $5$ divise $(y - 1)$.
Il existe donc un entier $k$ tel que $y - 1 = 5k$, d'où $y = 5k + 1$.
4. \textbf{Substitution :} On remplace $y - 1$ par $5k$ dans l'égalité $5(x - 1) = 3(y - 1)$ :
$$5(x - 1) = 3(5k) \implies x - 1 = 3k \implies x = 3k + 1$$
Les solutions sont donc les couples de la forme $(3k + 1, 5k + 1)$ avec $k \in \mathbb{Z}$.

\subsection{Exercice pratique avec Python}

\textbf{Objectif :} Implémenter l'algorithme d'Euclide étendu. 

Cet algorithme ne se contente pas de trouver le PGCD, il "remonte" les calculs pour trouver les coefficients $u$ et $v$ de l'identité de Bézout ($au + bv = \text{PGCD}$). C'est indispensable pour calculer des clés cryptographiques.

\begin{verbatim}
def euclide_etendu(a, b):
    """
    Retourne (pgcd, u, v) tels que a*u + b*v = pgcd
    basé sur l'algorithme d'Euclide étendu.
    """
    if b == 0:
        return a, 1, 0
    else:
        # Appel récursif
        d, u_prime, v_prime = euclide_etendu(b, a % b)
        
        # On reconstitue les coefficients u et v
        u = v_prime
        v = u_prime - (a // b) * v_prime
        
        return d, u, v

# --- Tests ---
a, b = 420, 135
pgcd, u, v = euclide_etendu(a, b)

print(f"Pour a={a} et b={b} :")
print(f"PGCD trouvé : {pgcd}")
print(f"Coefficients de Bézout : u={u}, v={v}")
print(f"Vérification : {a}*({u}) + {b}*({v}) = {a*u + b*v}")

# Test avec le problème d'introduction (récipients de 5L et 3L)
d, u, v = euclide_etendu(5, 3)
print("\nProblème des récipients (5L et 3L) :")
print(f"5*({u}) + 3*({v}) = {d}")
# Interprétation : Remplir 2 fois le récipient de 5L (10L), 
# et vider 3 fois son contenu dans celui de 3L (-9L) laisse 1L.
\end{verbatim}
\section*{Travaux Dirigés : Arithmétique dans $\mathbb{Z}$}
\addcontentsline{toc}{section}{Travaux Dirigés : Arithmétique dans $\mathbb{Z}$}

\vspace{0.5cm}

% --- Exercice 1 ---
\noindent\textbf{Exercice 1 : Restitution et définitions} \textit{(Taxonomie : Connaissance | Difficulté : Facile)}
\begin{enumerate}
    \item Énoncer avec précision le théorème de la division euclidienne dans $\mathbb{Z}$, en insistant sur la condition portant sur le reste.
    \item Soit $n \in \mathbb{N}^*$. Donner la définition mathématique exacte de la proposition : « $a$ et $b$ sont congrus modulo $n$ ».
    \item Rappeler l'énoncé du théorème de Gauss.
\end{enumerate}

\vspace{0.5cm}

% --- Exercice 2 ---
\noindent\textbf{Exercice 2 : Propriétés fondamentales} \textit{(Taxonomie : Compréhension | Difficulté : Facile)} \\
L'objectif de cet exercice est de démontrer des propriétés du cours à partir des définitions.
\begin{enumerate}
    \item Soient $a, b, c, d \in \mathbb{Z}$ et $n \in \mathbb{N}^*$. Démontrer que si $a \equiv b \pmod{n}$ et $c \equiv d \pmod{n}$, alors $a + c \equiv b + d \pmod{n}$.
    \item Démontrer que si $a \equiv b \pmod{n}$, alors pour tout entier naturel $k$, $a^k \equiv b^k \pmod{n}$.
\end{enumerate}

\vspace{0.5cm}

% --- Exercice 3 ---
\noindent\textbf{Exercice 3 : PGCD et littéral} \textit{(Taxonomie : Application | Difficulté : Moyenne)} \\
Soit $n$ un entier naturel. On considère les entiers $A = n$ et $B = n+1$.
\begin{enumerate}
    \item En utilisant la définition des diviseurs communs, démontrer que $\text{PGCD}(A, B) = 1$.
    \item En déduire, en appliquant l'identité de Bézout, que pour tout $n \in \mathbb{N}$, les nombres $2n+1$ et $n(n+1)$ sont premiers entre eux.
\end{enumerate}

\vspace{0.5cm}

% --- Exercice 4 ---
\noindent\textbf{Exercice 4 : Réduction de PGCD} \textit{(Taxonomie : Analyse | Difficulté : Moyenne)} \\
Soient $a$ et $b$ deux entiers relatifs non nuls. On pose $d = \text{PGCD}(a, b)$.
\begin{enumerate}
    \item Justifier l'existence de deux entiers $a'$ et $b'$ tels que $a = da'$ et $b = db'$.
    \item Analyser la relation entre $a'$ et $b'$ en démontrant rigoureusement que $\text{PGCD}(a', b') = 1$. \textit{(Indication : Utiliser le théorème de Bachet-Bézout)}.
\end{enumerate}

\vspace{0.5cm}

% --- Exercice 5 ---
\noindent\textbf{Exercice 5 : Conséquences du théorème de Gauss} \textit{(Taxonomie : Synthèse | Difficulté : Difficile)} \\
Soient $a, b$ et $c$ trois entiers relatifs non nuls. L'objectif est de démontrer une propriété de divisibilité simultanée.
\begin{enumerate}
    \item On suppose que $a$ divise $c$, que $b$ divise $c$, et que $\text{PGCD}(a, b) = 1$. Construire une démonstration prouvant que le produit $ab$ divise $c$. \textit{(Indication : Commencer par écrire $c = ak = bl$, puis utiliser le théorème de Gauss)}.
    \item Montrer par un contre-exemple que la condition $\text{PGCD}(a, b) = 1$ est strictement nécessaire.
\end{enumerate}

\vspace{0.5cm}

% --- Exercice 6 ---
\noindent\textbf{Exercice 6 : Évaluation de divisibilité} \textit{(Taxonomie : Évaluation | Difficulté : Difficile)} \\
On considère l'expression $E(n) = n(n+1)(2n+1)$ pour tout entier naturel $n$.
\begin{enumerate}
    \item Évaluer la divisibilité de $E(n)$ par $2$ en distinguant les cas de parité de $n$.
    \item Évaluer la divisibilité de $E(n)$ par $3$ en utilisant les congruences modulo $3$.
    \item En déduire, à l'aide du résultat de l'Exercice 5, que $E(n)$ est toujours un multiple de $6$.
\end{enumerate}

\vspace{0.5cm}

% --- Exercice 7 ---
\noindent\textbf{Exercice 7 : Somme et produit d'entiers premiers entre eux} \textit{(Taxonomie : Synthèse | Difficulté : Très Difficile)} \\
Soient $a$ et $b$ deux entiers relatifs non nuls tels que $\text{PGCD}(a, b) = 1$.
\begin{enumerate}
    \item Démontrer que $\text{PGCD}(a+b, a) = 1$ et que $\text{PGCD}(a+b, b) = 1$.
    \item En déduire rigoureusement que $\text{PGCD}(a+b, ab) = 1$.
    \item Cette propriété est-elle conservée pour la somme des carrés ? Démontrer que $\text{PGCD}(a^2+b^2, ab) = 1$.
\end{enumerate}

\vspace{0.5cm}

% --- Exercice 8 ---
\noindent\textbf{Exercice 8 : Construction du système des restes chinois (cas n=2)} \textit{(Taxonomie : Création | Difficulté : Très Difficile)} \\
Soient $p$ et $q$ deux entiers naturels non nuls et premiers entre eux ($\text{PGCD}(p, q) = 1$). On s'intéresse au système de congruences suivant, où $a$ et $b$ sont des entiers quelconques :
$$
(S) \begin{cases} 
x \equiv a \pmod{p} \\ 
x \equiv b \pmod{q} 
\end{cases}
$$
\begin{enumerate}
    \item Justifier l'existence de deux entiers $u$ et $v$ tels que $pu + qv = 1$.
    \item Créer une solution particulière $x_0$ au système $(S)$ en fonction de $p, q, u, v, a$ et $b$. Démontrer que cette valeur $x_0$ satisfait bien les deux équations du système.
    \item Formuler et démontrer la forme générale de toutes les solutions $x$ du système $(S)$. Sur quel modulo l'unicité de la solution est-elle garantie ?
\end{enumerate}
\section*{Travaux Pratiques (TP) : Implémentations Algorithmiques}
\addcontentsline{toc}{section}{Travaux Pratiques (TP) : Implémentations Algorithmiques}

\vspace{0.5cm}
\noindent L'objectif de cette séance de Travaux Pratiques est de traduire les théorèmes d'arithmétique dans l'ensemble $\mathbb{Z}$ en algorithmes fonctionnels. Les implémentations pourront être réalisées en langage Python.

\vspace{0.5cm}

% --- Exercice 1 ---
\noindent\textbf{Exercice 1 : La division euclidienne \textit{from scratch}} 
\begin{enumerate}
    \item \textbf{[Bloom : Compréhension | Difficulté : Facile]} \\
    Écrire une fonction \texttt{division\_naïve(a, b)} qui prend deux entiers naturels $a$ et $b$ ($b > 0$) et qui retourne le quotient $q$ et le reste $r$ en utilisant uniquement des soustractions successives (sans utiliser les opérateurs \texttt{//} et \texttt{\%}).
    \item \textbf{[Bloom : Application | Difficulté : Moyenne]} \\
    Adapter cette fonction pour qu'elle respecte le théorème mathématique de la division euclidienne dans $\mathbb{Z}$, c'est-à-dire en gérant correctement les cas où $a$ et/ou $b$ sont des entiers strictement négatifs (rappel : le reste $r$ doit toujours vérifier $0 \le r < |b|$).
\end{enumerate}

\vspace{0.5cm}

% --- Exercice 2 ---
\noindent\textbf{Exercice 2 : Congruences et tests de primalité} 
\begin{enumerate}
    \item \textbf{[Bloom : Application | Difficulté : Facile]} \\
    Écrire une fonction \texttt{est\_premier(n)} qui utilise l'opérateur modulo pour déterminer si un entier $n \ge 2$ est un nombre premier. L'algorithme devra s'arrêter intelligemment à $\lfloor\sqrt{n}\rfloor$ pour optimiser le calcul.
    \item \textbf{[Bloom : Analyse | Difficulté : Moyenne]} \\
    Implémenter une fonction \texttt{crible\_eratosthene(N)} qui retourne la liste de tous les nombres premiers inférieurs ou égaux à $N$, en appliquant le principe d'élimination des multiples.
\end{enumerate}

\vspace{0.5cm}

% --- Exercice 3 ---
\noindent\textbf{Exercice 3 : Les deux visages du PGCD} 
\begin{enumerate}
    \item \textbf{[Bloom : Application | Difficulté : Facile]} \\
    Implémenter une fonction \texttt{pgcd\_soustraction(a, b)} basée sur l'algorithme original d'Euclide utilisant les soustractions successives.
    \item \textbf{[Bloom : Application | Difficulté : Moyenne]} \\
    Implémenter une fonction \texttt{pgcd\_modulo(a, b)} utilisant les divisions euclidiennes successives (opérateur \texttt{\%}). 
    \item \textbf{[Bloom : Évaluation | Difficulté : Moyenne]} \\
    Écrire un script qui compare le temps d'exécution de ces deux fonctions pour $a = 123456789$ et $b = 123$. Conclure sur la complexité algorithmique.
\end{enumerate}

\vspace{0.5cm}

% --- Exercice 4 ---
\noindent\textbf{Exercice 4 : Calcul du PPCM optimisé} 
\begin{enumerate}
    \item \textbf{[Bloom : Compréhension | Difficulté : Facile]} \\
    Déduire du théorème vu en cours une fonction \texttt{ppcm(a, b)} qui calcule le Plus Petit Commun Multiple de deux entiers en faisant appel à la fonction \texttt{pgcd\_modulo(a, b)}.
    \item \textbf{[Bloom : Synthèse | Difficulté : Moyenne]} \\
    Écrire une fonction \texttt{ppcm\_liste(L)} qui prend en paramètre une liste de $n$ entiers relatifs et qui retourne le PPCM global de tous les éléments de cette liste.
\end{enumerate}

\vspace{0.5cm}

% --- Exercice 5 ---
\noindent\textbf{Exercice 5 : L'Algorithme d'Euclide Étendu} 
\begin{enumerate}
    \item \textbf{[Bloom : Analyse | Difficulté : Difficile]} \\
    Implémenter la fonction \texttt{euclide\_etendu(a, b)} qui retourne un triplet $(d, u, v)$ tel que $d = \text{PGCD}(a, b)$ et $au + bv = d$. Vous pouvez utiliser une approche récursive ou itérative au choix.
    \item \textbf{[Bloom : Évaluation | Difficulté : Moyenne]} \\
    Tester la fonction avec $a = 240$ et $b = 46$, puis afficher la chaîne de caractères vérifiant l'équation de Bézout correspondante.
\end{enumerate}

\vspace{0.5cm}

% --- Exercice 6 ---
\noindent\textbf{Exercice 6 : Résolution algorithmique d'équations diophantiennes} 
\begin{enumerate}
    \item \textbf{[Bloom : Synthèse | Difficulté : Difficile]} \\
    Écrire une fonction \texttt{resoudre\_diophantienne(a, b, c)} qui cherche à résoudre l'équation $ax + by = c$. La fonction devra vérifier si une solution existe grâce au théorème de Bachet-Bézout, et si oui, utiliser la fonction \texttt{euclide\_etendu} pour retourner une solution particulière $(x_0, y_0)$.
    \item \textbf{[Bloom : Création | Difficulté : Difficile]} \\
    Enrichir la fonction pour qu'elle prenne en compte deux bornes $X_{\text{min}}$ et $X_{\text{max}}$, et qu'elle retourne la liste de toutes les solutions $(x, y)$ telles que $X_{\text{min}} \le x \le X_{\text{max}}$.
\end{enumerate}

\vspace{0.5cm}

% --- Exercice 7 ---
\noindent\textbf{Exercice 7 : Application Cryptographique (Chiffrement Affine)} 
\begin{enumerate}
    \item \textbf{[Bloom : Application | Difficulté : Moyenne]} \\
    Le chiffrement affine transforme une lettre en lui associant un rang $x$ (de 0 pour A à 25 pour Z), puis calcule un nouveau rang $y \equiv ax + b \pmod{26}$. Écrire une fonction \texttt{chiffrement\_affine(texte, a, b)} qui retourne le texte chiffré. Que doit vérifier la clé $a$ pour que le déchiffrement soit possible ?
    \item \textbf{[Bloom : Synthèse | Difficulté : Difficile]} \\
    Pour déchiffrer, il faut trouver l'inverse modulaire de $a$ modulo 26, c'est-à-dire un entier $a'$ tel que $a \times a' \equiv 1 \pmod{26}$. Utiliser l'algorithme d'Euclide étendu pour écrire une fonction \texttt{inverse\_modulaire(a, n)} et en déduire la fonction \texttt{dechiffrement\_affine(texte\_chiffre, a, b)}.
\end{enumerate}

\vspace{0.5cm}

% --- Exercice 8 ---
\noindent\textbf{Exercice 8 : Introduction à RSA (Génération de clés)} 
\begin{enumerate}
    \item \textbf{[Bloom : Création | Difficulté : Très Difficile]} \\
    L'algorithme de cryptographie asymétrique RSA repose fortement sur l'arithmétique dans $\mathbb{Z}$. Écrire une fonction globale \texttt{generer\_cles\_rsa(p, q)} où $p$ et $q$ sont deux nombres premiers distincts donnés en paramètre.
    L'algorithme devra :
    \begin{itemize}
        \item Calculer le module $n = p \times q$.
        \item Calculer l'indicatrice d'Euler $\phi(n) = (p-1)(q-1)$.
        \item Choisir un entier $e$ tel que $1 < e < \phi(n)$ et $\text{PGCD}(e, \phi(n)) = 1$.
        \item Calculer $d$, l'inverse modulaire de $e$ modulo $\phi(n)$ (à l'aide de l'exercice 7).
        \item Retourner la clé publique $(n, e)$ et la clé privée $(n, d)$.
    \end{itemize}
\end{enumerate}
\section*{Aller plus loin : L'arithmétique au cœur de la Data Science}
\addcontentsline{toc}{section}{Aller plus loin : L'arithmétique au cœur de la Data Science}

L'arithmétique dans $\mathbb{Z}$, bien qu'étudiée depuis l'Antiquité, est aujourd'hui l'un des piliers de l'ingénierie des données (Data Science) et de l'intelligence artificielle. Au-delà de la cryptographie (RSA) que nous avons évoquée, les concepts de division euclidienne et de congruence sont indispensables pour optimiser la mémoire et les temps de calcul dans les algorithmes d'apprentissage automatique (Machine Learning).

\subsection*{Le défi des données textuelles en NLP}

Dans le domaine du traitement du langage naturel (NLP), les algorithmes ne comprennent pas les mots : ils ont besoin de nombres. Pour entraîner un modèle de classification de textes ou analyser des séquences de mots (ce qu'on appelle des \textit{N-grammes}), on transforme généralement le vocabulaire en un immense tableau (un dictionnaire). 

Cependant, sur un corpus de texte réel (comme l'ensemble des articles Wikipédia), le nombre de mots uniques et de combinaisons de $N$-grammes peut atteindre des dizaines de millions. Créer une matrice où chaque colonne représente un $N$-gramme spécifique demande une quantité de mémoire RAM gigantesque, souvent indisponible sur des machines standards. Comment alors représenter ces données spatialement sans saturer la mémoire ?

\subsection*{La solution : Le « Feature Hashing » (L'astuce du hachage)}

C'est ici qu'intervient l'arithmétique modulaire à travers une technique appelée le \textbf{Feature Hashing} (ou Hashing Trick). Plutôt que de maintenir un dictionnaire exact, on utilise une fonction mathématique (une fonction de hachage) qui transforme chaque chaîne de caractères en un grand nombre entier abstrait.

Ensuite, pour forcer ces grands nombres à rentrer dans un espace mémoire de taille fixe et contrôlée (par exemple $M$ colonnes), on applique une simple congruence modulo $M$ :
$$ \text{Index du vecteur} = \text{Fonction\_Hachage}(\text{mot}) \pmod M $$



\textbf{Exemple conceptuel :}
Supposons que nous fixions la taille de notre mémoire à $M = 1000$ dimensions. 
Si la fonction de hachage convertit le bi-gramme \textit{"intelligence artificielle"} en l'entier $4598273$, sa position dans notre vecteur mémoire sera simplement le reste de sa division euclidienne par $1000$ :
$$ 4598273 \equiv 273 \pmod{1000} $$
Le modèle stockera l'information de ce bi-gramme à l'index 273.

\subsection*{Gestion des collisions et propriétés statistiques}

Que se passe-t-il si un autre mot possède également un reste de 273 (ce qu'on appelle une collision) ? Mathématiquement, le modèle va additionner les fréquences de ces deux mots au même endroit. 

Bien que cela semble être une perte d'information (deux mots différents mélangés), les mathématiques statistiques démontrent que si $M$ est suffisamment grand (et souvent choisi comme un nombre premier pour de meilleures propriétés de répartition arithmétique), l'impact de ces collisions sur les performances du modèle de Machine Learning est négligeable. 

Ainsi, grâce à la théorie des congruences, des modèles très complexes peuvent s'entraîner sur des flux de données continus et infinis avec une empreinte mémoire fixe et strictement maîtrisée en temps constant $\mathcal{O}(1)$.

\vspace{0.5cm}
\begin{center}
\textit{Pour aller plus loin sur ce sujet, vous pouvez vous documenter sur l'algorithme \textbf{MurmurHash}, très utilisé dans les bibliothèques de Machine Learning comme Scikit-Learn (via \texttt{FeatureHasher}), ou sur la manière dont les architectures de plongement lexical (Word Embeddings) comme \textbf{GloVe} gèrent les matrices de co-occurrence massives.}
\end{center}
\section*{Projet d'Application : Le « Hashing Trick » en Analyse de Données}
\addcontentsline{toc}{section}{Projet d'Application : Le « Hashing Trick » en Analyse de Données}

\subsection*{Contexte du Projet}
Imaginez que vous avez mis en place un système de \textit{web scraping} pour récupérer en direct les commentaires textuels de milliers de matchs de football (ex: « Tir cadré de l'attaquant », « Faute dangereuse au milieu du terrain », etc.). Votre objectif est d'entraîner un modèle de Machine Learning capable d'analyser ces flux de textes pour calculer des statistiques en temps réel ou prédire l'issue d'une action.

Le problème : le vocabulaire total généré par ces milliers de matchs est gigantesque. Construire un dictionnaire classique en mémoire vive (RAM) pour associer chaque mot ou N-gramme à un index unique va saturer vos serveurs. 

Vous allez donc utiliser l'arithmétique dans $\mathbb{Z}$, et plus précisément la division euclidienne, pour implémenter de zéro la technique du \textbf{Feature Hashing} étudiée dans la section précédente.

\subsection*{Objectifs}
\begin{enumerate}
    \item Implémenter une vectorisation de texte classique (avec dictionnaire) pour en montrer les limites.
    \item Implémenter une fonction de \textit{Feature Hashing} basée sur l'opération modulo.
    \item Analyser mathématiquement le taux de collision en fonction de la taille de l'espace mémoire $M$.
\end{enumerate}

\subsection*{Cahier des charges et Travail demandé}

\subsubsection*{Étape 1 : Simulation des données textuelles}
Pour simuler notre flux de données, créez un script Python qui génère une liste de phrases ou récupérez un corpus de texte volumineux (par exemple, un fichier texte contenant un grand nombre de commentaires sportifs ou d'articles).
Développez une fonction \texttt{extraire\_mots(texte)} qui nettoie le texte (minuscules, ponctuation) et retourne une liste de mots.

\subsubsection*{Étape 2 : L'approche classique (Dictionnaire exact)}
Écrivez une fonction \texttt{vectoriseur\_classique(corpus)} qui :
\begin{itemize}
    \item Parcourt tous les mots du corpus.
    \item Assigne un identifiant unique (un entier de $0$ à $V-1$, où $V$ est la taille du vocabulaire) à chaque nouveau mot rencontré en utilisant un dictionnaire Python (\texttt{dict}).
    \item \textbf{Question :} Quelle est la complexité spatiale de cette approche si le flux de mots est infini ?
\end{itemize}

\subsubsection*{Étape 3 : L'approche arithmétique (Feature Hashing)}
Ici, nous n'utilisons plus de dictionnaire pour mémoriser les mots. Nous fixons une taille de mémoire maximale, par exemple $M = 1024$ colonnes.

Écrivez une fonction \texttt{vectoriseur\_hachage(corpus, M)} qui :
\begin{itemize}
    \item Utilise une fonction de hachage standard pour convertir chaque mot en un entier abstrait. Vous pouvez utiliser la fonction native \texttt{hash()} de Python ou importer \texttt{hashlib}.
    \item Calcule l'index du mot dans le vecteur mémoire en utilisant la congruence : 
    $$ \text{Index} \equiv \text{hash}(\text{mot}) \pmod M $$
    \item Construit un vecteur de taille $M$ comptant la fréquence d'apparition des index.
\end{itemize}

\subsubsection*{Étape 4 : Analyse arithmétique des collisions}
Une collision survient lorsque deux mots différents $m_1$ et $m_2$ vérifient :
$$ \text{hash}(m_1) \equiv \text{hash}(m_2) \pmod M $$
\begin{enumerate}
    \item Modifiez votre fonction de hachage pour qu'elle compte le nombre exact de mots différents qui "tombent" sur le même index.
    \item Testez votre algorithme avec $M = 100$, $M = 1009$ (un nombre premier) et $M = 10000$.
    \item Tracez un graphique (avec \texttt{matplotlib}) représentant le nombre de collisions en fonction de la taille $M$. 
    \item \textbf{Analyse :} Pourquoi recommande-t-on souvent d'utiliser un nombre premier pour la dimension $M$ en Data Science ? (Faites le lien avec les diviseurs communs vus dans le Chapitre III).
\end{enumerate}

\subsubsection*{Bonus (Pour aller très loin)}
Remplacez les mots simples par des bi-grammes (paires de mots adjacents). Comparez l'explosion de la mémoire RAM avec l'approche classique par rapport à la stabilité absolue de la mémoire avec votre approche par modulo.

\input{ch4}
\chapter{Polynômes}

Les polynômes constituent l'outil de calcul algébrique le plus polyvalent : interpolation et approximation numérique, traitement du signal, codes correcteurs (Reed-Solomon), cryptographie sur courbes elliptiques, méthodes à noyaux en apprentissage automatique. Ce chapitre développe leur arithmétique : structure d'anneau, racines, division euclidienne, factorisation et irréductibilité. Il prolonge directement le chapitre III (arithmétique dans $\mathbb{Z}$) et le chapitre IV (anneaux et corps).

Dans tout ce chapitre, $K$ désigne un corps commutatif, le plus souvent $\mathbb{R}$ ou $\mathbb{C}$.

\section{Anneaux $\mathbb{R}[X]$ et $\mathbb{C}[X]$}

\subsection{Motivation}
Un polynôme « scolaire » comme $P(X) = X^2 - 3X + 2$ peut être vu de deux manières :
\begin{itemize}
    \item Comme un \emph{objet formel} (la suite des coefficients $(2, -3, 1, 0, 0, \dots)$) que l'on additionne et multiplie selon des règles algébriques.
    \item Comme une \emph{fonction} $x \mapsto x^2 - 3x + 2$ que l'on évalue en chaque point.
\end{itemize}
Sur $\mathbb{R}$ ou $\mathbb{C}$, les deux points de vue sont équivalents : un polynôme est entièrement déterminé par sa fonction. Mais sur un corps fini, ce n'est plus vrai. Cette section pose les bases formelles : $K[X]$ comme \emph{anneau abstrait}, indépendamment de l'évaluation. La distinction sera levée à la section suivante (§2).

\subsection{Approche par problème : polynôme nul vs fonction nulle}
\textbf{Problème :} Considérons sur le corps $\mathbb{F}_2 = \mathbb{Z}/2\mathbb{Z}$ le polynôme $P(X) = X^2 + X$. Calculer $P(\bar{0})$ et $P(\bar{1})$. Que peut-on dire ? Le polynôme $X^2 + X$ est-il « nul » ?

\textbf{Résolution :} $P(\bar{0}) = \bar{0}$ et $P(\bar{1}) = \bar{1}^2 + \bar{1} = \bar{1} + \bar{1} = \bar{0}$. La fonction $x \mapsto P(x)$ est donc identiquement nulle sur $\mathbb{F}_2$. Pourtant le polynôme $X^2 + X$ a des coefficients non tous nuls : il \emph{n'est pas} le polynôme nul. Conclusion : sur $\mathbb{F}_2$, deux polynômes distincts peuvent définir la même fonction. Sur $\mathbb{R}$ ou $\mathbb{C}$, ce phénomène n'arrive pas. Il faut donc distinguer rigoureusement « polynôme » (objet algébrique formel) et « fonction polynomiale ».

\subsection{Construction formelle de $K[X]$}

\begin{definition}[Polynôme à coefficients dans $K$]
Un \textbf{polynôme} à coefficients dans $K$ est une suite $(a_n)_{n \in \mathbb{N}}$ d'éléments de $K$ \emph{presque tous nuls} (i.e.\ tous nuls à partir d'un certain rang). On note l'ensemble de ces suites $K[X]$.

Si $(a_n) \in K[X]$, on l'écrit :
\[ P = \sum_{n=0}^{+\infty} a_n X^n = a_0 + a_1 X + a_2 X^2 + \cdots \]
en utilisant le symbole $X$ comme \emph{indéterminée}. Les $a_n$ sont les \textbf{coefficients} de $P$. Le \textbf{polynôme nul}, noté $0$, est la suite identiquement nulle.
\end{definition}

\begin{definition}[{Opérations sur $K[X]$}]
Soient $P = \sum a_n X^n$ et $Q = \sum b_n X^n$ deux polynômes de $K[X]$, et $\lambda \in K$.
\begin{align*}
P + Q &= \sum_{n} (a_n + b_n) X^n, \\
\lambda P &= \sum_{n} (\lambda a_n) X^n, \\
P \times Q &= \sum_{n} c_n X^n \quad \text{où} \quad c_n = \sum_{k=0}^{n} a_k\, b_{n-k}.
\end{align*}
La formule du produit est appelée \emph{produit de Cauchy}. Elle correspond exactement à la multiplication usuelle des polynômes.
\end{definition}

\begin{theoreme}[{$K[X]$ est un anneau commutatif unitaire}]
Muni de $+$ et $\times$ ci-dessus, $K[X]$ est un \emph{anneau commutatif unitaire}. Le neutre additif est le polynôme nul, le neutre multiplicatif est le polynôme constant $1$.
\end{theoreme}

\subsection{Degré et valuation}

\begin{definition}[Degré]
Soit $P = \sum a_n X^n \in K[X]$ non nul. Le \textbf{degré} de $P$, noté $\deg P$, est le plus grand entier $n$ tel que $a_n \neq 0$. On pose par convention $\deg 0 = -\infty$.

Le coefficient $a_{\deg P}$ est le \textbf{coefficient dominant} de $P$. Si ce coefficient vaut $1$, le polynôme est dit \textbf{unitaire} (ou monique).
\end{definition}

\begin{definition}[Valuation]
La \textbf{valuation} de $P \neq 0$, notée $v(P)$, est le plus petit entier $n$ tel que $a_n \neq 0$. Par convention, $v(0) = +\infty$.
\end{definition}

\begin{proposition}[Propriétés du degré]
Pour tous $P, Q \in K[X]$ :
\begin{enumerate}
    \item $\deg(P + Q) \le \max(\deg P, \deg Q)$, avec égalité si $\deg P \neq \deg Q$.
    \item $\deg(P \times Q) = \deg P + \deg Q$.
    \item $\deg(\lambda P) = \deg P$ pour tout $\lambda \in K^*$.
\end{enumerate}
(Conventions : $-\infty + n = -\infty$, $\max(-\infty, n) = n$.)
\end{proposition}

\begin{proof}
(1) Le coefficient de $X^n$ dans $P + Q$ est $a_n + b_n$, qui est nul si $n > \max(\deg P, \deg Q)$. Si $\deg P > \deg Q$, le coefficient de $X^{\deg P}$ dans $P + Q$ vaut $a_{\deg P} \neq 0$.

(2) Le coefficient dominant de $PQ$ est $a_{\deg P} \cdot b_{\deg Q}$, non nul puisque $K$ est un corps (donc intègre).

(3) Conséquence immédiate de (2) avec $Q = \lambda$ (constante non nulle).
\end{proof}

\begin{corollaire}[{Intégrité de $K[X]$}]
Pour tout corps $K$, l'anneau $K[X]$ est \emph{intègre}. En particulier, $\mathbb{R}[X]$ et $\mathbb{C}[X]$ sont des anneaux commutatifs unitaires intègres.
\end{corollaire}

\begin{proof}
Si $P \neq 0$ et $Q \neq 0$, alors $\deg(PQ) = \deg P + \deg Q \ge 0$, donc $PQ \neq 0$.
\end{proof}

\subsection{Inversibles de $K[X]$}

\begin{proposition}[Inversibles]
Les éléments inversibles de $K[X]$ pour la multiplication sont exactement les \emph{constantes non nulles} :
\[ K[X]^\times = K^* = K \setminus \{0\}. \]
\end{proposition}

\begin{proof}
Si $P \in K^*$, son inverse dans $K$ est aussi son inverse dans $K[X]$.
Réciproquement, si $PQ = 1$, alors $\deg P + \deg Q = \deg 1 = 0$. Comme les degrés sont $\ge 0$, on a $\deg P = \deg Q = 0$ : $P$ et $Q$ sont des constantes non nulles.
\end{proof}

\begin{remarque}
Ce résultat est analogue à $\mathbb{Z}^\times = \{-1, 1\}$ : un anneau intègre n'est pas un corps en général. C'est précisément ce qui rend l'étude arithmétique de $K[X]$ riche (divisibilité, irréductibilité, factorisation), à l'image de l'arithmétique dans $\mathbb{Z}$.
\end{remarque}

\subsection{Cas particuliers : $\mathbb{R}[X]$ et $\mathbb{C}[X]$}

Les corps $\mathbb{R}$ et $\mathbb{C}$ sont les terrains de jeu privilégiés du chapitre.

\begin{itemize}
    \item $\mathbb{R}[X]$ et $\mathbb{C}[X]$ sont tous deux des anneaux commutatifs unitaires intègres.
    \item L'inclusion $\mathbb{R} \subset \mathbb{C}$ induit une inclusion $\mathbb{R}[X] \subset \mathbb{C}[X]$ (un polynôme à coefficients réels est en particulier à coefficients complexes).
    \item La conjugaison complexe se prolonge en un automorphisme de $\mathbb{C}[X]$ : si $P = \sum a_k X^k \in \mathbb{C}[X]$, on note $\bar{P} = \sum \bar{a_k} X^k$. On a $P \in \mathbb{R}[X] \iff \bar{P} = P$.
    \item Le théorème fondamental de l'algèbre (admis ici, prouvé en analyse complexe) affirme que tout polynôme non constant de $\mathbb{C}[X]$ admet au moins une racine dans $\mathbb{C}$. Cette propriété fera de $\mathbb{C}$ un corps \emph{algébriquement clos}, et structurera toute la décomposition étudiée au §4.
\end{itemize}

\subsection{Composition de polynômes}

\begin{definition}[Composition]
Soient $P = \sum a_k X^k$ et $Q \in K[X]$. La \textbf{composée} de $P$ par $Q$ est :
\[ P \circ Q = P(Q) = \sum_{k=0}^{\deg P} a_k\, Q^k \in K[X]. \]
\end{definition}

\begin{proposition}
Si $\deg P, \deg Q \ge 1$, alors $\deg(P \circ Q) = \deg P \cdot \deg Q$.
\end{proposition}

\subsection{Exemples résolus (Difficulté ascendante)}

\textbf{Exemple 1 (Calcul direct --- Facile).}
Soient $P = 2X^3 - X + 5$ et $Q = X^2 + 3$ dans $\mathbb{R}[X]$. Calculer $P + Q$, $P \times Q$ et leurs degrés.

\textit{Correction :}
$P + Q = 2X^3 + X^2 - X + 8$, de degré $3$.
$P \times Q = (2X^3)(X^2 + 3) + (-X)(X^2 + 3) + 5(X^2 + 3) = 2X^5 + 6X^3 - X^3 - 3X + 5X^2 + 15 = 2X^5 + 5X^3 + 5X^2 - 3X + 15$, de degré $5 = 3 + 2$. \checkmark

\medskip

\textbf{Exemple 2 (Conjugaison --- Intermédiaire).}
Soit $P = X^3 + (2 - i)X + 4i \in \mathbb{C}[X]$. Calculer $\bar{P}$, puis $P + \bar{P}$ et $P \cdot \bar{P}$. Lequel des deux résultats appartient à $\mathbb{R}[X]$ ?

\textit{Correction :}
$\bar{P} = X^3 + (2 + i)X - 4i$.
$P + \bar{P} = 2X^3 + 4X$ : à coefficients réels, donc dans $\mathbb{R}[X]$. \checkmark
$P \cdot \bar{P}$ a aussi tous ses coefficients réels (on multiplie un polynôme par son conjugué). En général, $P + \bar{P}$, $P \cdot \bar{P}$ et $i(P - \bar{P})$ sont des polynômes réels, peu importe $P \in \mathbb{C}[X]$.

\medskip

\textbf{Exemple 3 (Identification de coefficients --- Difficile).}
Trouver tous les polynômes $P \in \mathbb{R}[X]$ tels que $P(X+1) - P(X) = X$.

\textit{Correction :} si $\deg P = n$, alors $\deg(P(X+1) - P(X)) = n - 1$ (le coefficient dominant disparaît par soustraction). Comme le membre de droite a degré $1$, on a $n = 2$. Posons $P = aX^2 + bX + c$.
$P(X+1) - P(X) = a(X+1)^2 + b(X+1) + c - aX^2 - bX - c = a(2X + 1) + b = 2aX + a + b$.
Identification avec $X$ : $2a = 1$ et $a + b = 0$, soit $a = 1/2$ et $b = -1/2$. La constante $c$ reste libre.
Donc $P(X) = \frac{X^2 - X}{2} + c = \frac{X(X-1)}{2} + c$.

(On reconnaît $\binom{X}{2}$, lié au calcul de $\sum_{k=0}^{n-1} k$.)

\subsection{Exercices pratiques avec Python}
\texttt{numpy.polynomial.Polynomial} et \texttt{sympy} permettent de manipuler symboliquement et numériquement les polynômes.

\begin{lstlisting}
import numpy as np
from numpy.polynomial import Polynomial
import sympy as sp

# 1. Numpy : polynomes a coefficients reels (ordre croissant des coefficients)
P = Polynomial([5, -1, 0, 2])   # 5 - X + 2 X^3
Q = Polynomial([3, 0, 1])       # 3 + X^2
print(f"P = {P}")
print(f"Q = {Q}")
print(f"deg(P) = {P.degree()}, deg(Q) = {Q.degree()}")
print(f"P + Q = {P + Q}")
print(f"P * Q = {P * Q}    (degre {(P*Q).degree()})")

# 2. Sympy : calcul symbolique exact
X = sp.symbols('X')
P_s = 2*X**3 - X + 5
Q_s = X**2 + 3
print(f"\nP * Q (sympy) = {sp.expand(P_s * Q_s)}")

# 3. Resolution de l'exemple 3 : trouver P tel que P(X+1) - P(X) = X
P_inconnu = sum(sp.Symbol(f'a{i}') * X**i for i in range(3))
eq = sp.expand(P_inconnu.subs(X, X + 1) - P_inconnu - X)
sol = sp.solve(sp.Poly(eq, X).all_coeffs(), [sp.Symbol('a0'), sp.Symbol('a1'), sp.Symbol('a2')])
print(f"\nSolution P : {sol}  (a0 reste libre)")
\end{lstlisting}

\subsection{Exercices à faire}
\begin{enumerate}
    \item \textbf{Exercice 1.} Soit $P = 3X^4 - 2X^2 + X - 1$ et $Q = X^3 + X$ dans $\mathbb{Q}[X]$. Calculer $P + Q$, $P - Q$, $PQ$ et leurs degrés.
    \item \textbf{Exercice 2.} Démontrer que $\deg(P \circ Q) = \deg P \cdot \deg Q$ lorsque $P, Q$ sont non constants. Que se passe-t-il si $Q$ est une constante ?
    \item \textbf{Exercice 3.} Trouver tous les $P \in \mathbb{R}[X]$ vérifiant $P(X^2) = (X^2 + 1) P(X)$. (Indication : commencer par étudier le degré, puis le coefficient dominant.)
    \item \textbf{Exercice 4.} Combien y a-t-il de polynômes unitaires de degré $3$ dans $\mathbb{F}_2[X]$ ? Les énumérer.
    \item \textbf{Exercice 5 (synthèse).} Soit $P \in \mathbb{C}[X]$. Montrer les équivalences :
    \begin{itemize}
        \item $P \in \mathbb{R}[X] \iff \bar{P} = P$ ;
        \item $P + \bar{P}$, $P \cdot \bar{P}$ et $i(P - \bar{P})$ sont toujours dans $\mathbb{R}[X]$.
    \end{itemize}
\end{enumerate}

\section{Fonctions polynomiales et racines}

\subsection{Motivation}
La résolution d'équations polynomiales est l'un des plus vieux problèmes des mathématiques. De la formule du second degré apprise au lycée à la résolution numérique en analyse, en passant par la cryptographie sur courbes elliptiques (où les opérations reposent sur des racines de polynômes de degré 3) et les codes correcteurs d'erreurs (Reed-Solomon détecte les erreurs en évaluant un polynôme), \emph{les racines structurent toute l'algèbre du calcul}. Cette section formalise le lien entre un polynôme (objet abstrait) et la fonction qu'il définit, et caractérise ses racines par leur multiplicité.

\subsection{Approche par problème : reconstruire un polynôme à partir de ses racines}
\textbf{Problème :} Un polynôme $P \in \mathbb{R}[X]$ unitaire de degré $5$ admet pour racines $1$ (simple), $2$ (double) et $1 + i$ (simple). Reconstruire $P$.

\textbf{Résolution :} Comme $P$ est à coefficients réels, $\overline{1 + i} = 1 - i$ est aussi racine, et de même multiplicité que $1 + i$ (théorème ci-dessous). On a donc $5$ racines comptées avec multiplicité : $1, 2, 2, 1+i, 1-i$. Comme $\deg P = 5$ et que $P$ est unitaire :
\[ P = (X - 1)(X - 2)^2 (X - (1+i))(X - (1-i)) = (X - 1)(X - 2)^2 (X^2 - 2X + 2). \]
On vérifie que $(X - (1+i))(X - (1-i)) = X^2 - 2X + (1+i)(1-i) = X^2 - 2X + 2$, à coefficients réels. \checkmark

\subsection{Évaluation et fonction polynomiale}

\begin{definition}[Évaluation, fonction polynomiale]
Soit $P = \sum_{k=0}^{n} a_k X^k \in K[X]$ et $a \in K$. L'\textbf{évaluation} de $P$ en $a$ est
\[ P(a) = \sum_{k=0}^{n} a_k\, a^k \in K. \]
La \textbf{fonction polynomiale} associée à $P$ est l'application $\widetilde{P} : K \to K,\ x \mapsto P(x)$.
\end{definition}

\begin{proposition}[Morphisme d'évaluation]
Pour tout $a \in K$, l'application $\mathrm{ev}_a : K[X] \to K,\ P \mapsto P(a)$ est un morphisme d'anneaux. Autrement dit :
\[ (P + Q)(a) = P(a) + Q(a), \qquad (P \cdot Q)(a) = P(a) \cdot Q(a), \qquad 1(a) = 1. \]
\end{proposition}

\subsection{Racines et factorisation par $(X - a)$}

\begin{definition}[Racine]
Un élément $a \in K$ est une \textbf{racine} (ou \emph{zéro}) de $P \in K[X]$ si $P(a) = 0$.
\end{definition}

\begin{theoreme}[Caractérisation par factorisation]
Soit $P \in K[X]$ et $a \in K$. Alors $a$ est racine de $P$ si et seulement si $(X - a)$ divise $P$ dans $K[X]$, c'est-à-dire qu'il existe $Q \in K[X]$ tel que $P = (X - a) Q$.
\end{theoreme}

\begin{proof}
($\Leftarrow$) Si $P = (X - a) Q$, alors $P(a) = 0 \cdot Q(a) = 0$.

($\Rightarrow$) On utilise l'identité $X^k - a^k = (X - a)(X^{k-1} + a X^{k-2} + \cdots + a^{k-1})$. Pour $P = \sum a_k X^k$, on a
\[ P(X) - P(a) = \sum_{k} a_k (X^k - a^k) = (X - a) \cdot \underbrace{\sum_{k} a_k (X^{k-1} + a X^{k-2} + \cdots + a^{k-1})}_{= Q(X)}. \]
Si $P(a) = 0$, ceci s'écrit $P = (X - a) Q$.
\end{proof}

\subsection{Multiplicité d'une racine}

\begin{definition}[Multiplicité]
Soit $a$ une racine de $P \in K[X]$, $P \neq 0$. La \textbf{multiplicité} de $a$ dans $P$ est l'entier $m \ge 1$ tel que :
\[ (X - a)^m \mid P \quad \text{et} \quad (X - a)^{m+1} \nmid P. \]
On la note $\mathrm{mult}_a(P)$. Si $m = 1$, on dit que $a$ est une racine \emph{simple}, si $m = 2$ \emph{double}, et plus généralement \emph{de multiplicité $m$}.

Si $a$ n'est pas racine, on pose $\mathrm{mult}_a(P) = 0$.
\end{definition}

\subsection{Polynôme dérivé}

\begin{definition}[Dérivée formelle]
Soit $P = \sum_{k=0}^{n} a_k X^k \in K[X]$. Le \textbf{polynôme dérivé} de $P$ est
\[ P' = \sum_{k=1}^{n} k\, a_k X^{k-1} \in K[X]. \]
On définit récursivement les dérivées d'ordre supérieur : $P^{(0)} = P$ et $P^{(m+1)} = (P^{(m)})'$.
\end{definition}

\begin{proposition}[Règles de dérivation]
La dérivation $P \mapsto P'$ est $K$-linéaire et vérifie la règle de Leibniz :
\[ (P + Q)' = P' + Q', \quad (\lambda P)' = \lambda P', \quad (PQ)' = P' Q + P Q'. \]
On en déduit la formule de la composée : $(P \circ Q)' = (P' \circ Q) \cdot Q'$.
Si $\mathrm{car}(K) = 0$ et $\deg P \ge 1$, alors $\deg P' = \deg P - 1$.
\end{proposition}

\begin{theoreme}[Caractérisation de la multiplicité par les dérivées]
Soit $K$ un corps de caractéristique $0$ (en particulier $\mathbb{Q}, \mathbb{R}, \mathbb{C}$). Soit $P \in K[X]$ non nul et $a \in K$. La multiplicité de $a$ dans $P$ est $m$ si et seulement si :
\[ P(a) = P'(a) = \cdots = P^{(m-1)}(a) = 0 \quad \text{et} \quad P^{(m)}(a) \neq 0. \]
\end{theoreme}

\begin{proof}
Écrivons $P = (X - a)^m\, Q$ avec $Q(a) \neq 0$. Par la règle de Leibniz, on calcule successivement les dérivées : pour $k < m$, $P^{(k)}(a) = 0$, et pour $k = m$, $P^{(m)}(a) = m!\, Q(a) \neq 0$ (en caractéristique $0$).
\end{proof}

\begin{remarque}
Cette caractérisation est \emph{fausse} en caractéristique non nulle. Sur $\mathbb{F}_p[X]$, le polynôme $(X - a)^p$ a pour dérivée $p(X-a)^{p-1} = 0$, ce qui invaliderait la formule. Dans nos applications usuelles ($\mathbb{R}, \mathbb{C}$), aucun problème.
\end{remarque}

\subsection{Borne sur le nombre de racines}

\begin{theoreme}[Nombre de racines dans un anneau intègre]
Soit $K$ un corps (ou plus généralement un anneau intègre) et $P \in K[X]$ non nul. Le nombre de racines de $P$ dans $K$, comptées avec multiplicité, est au plus $\deg P$.
\end{theoreme}

\begin{proof}
Par récurrence sur le degré. Si $\deg P = 0$, $P$ est une constante non nulle, sans racine. Pour $\deg P \ge 1$ : si $P$ a une racine $a$, on factorise $P = (X - a)^m Q$ avec $Q(a) \neq 0$, où $m \ge 1$ est la multiplicité. Alors $\deg Q = \deg P - m$. Par récurrence, $Q$ a au plus $\deg P - m$ racines, donc $P$ en a au plus $m + (\deg P - m) = \deg P$.
\end{proof}

\begin{corollaire}[Polynôme nul à beaucoup de racines]
Si $P \in K[X]$ admet strictement plus de $\deg P$ racines distinctes (ou comptées avec multiplicité), alors $P = 0$.
\end{corollaire}

\subsection{Identification polynôme-fonction}

\begin{theoreme}[Polynôme $\leftrightarrow$ fonction sur un corps infini]
Si $K$ est un corps infini (par exemple $\mathbb{Q}, \mathbb{R}, \mathbb{C}$), l'application $P \mapsto \widetilde{P}$ est injective. Autrement dit :
\[ \widetilde{P} = \widetilde{Q} \iff P = Q. \]
\end{theoreme}

\begin{proof}
Si $\widetilde{P} = \widetilde{Q}$, alors $\widetilde{P - Q} = 0$, donc $P - Q$ admet une infinité de racines (tout élément de $K$). Par le corollaire précédent, $P - Q = 0$.
\end{proof}

\begin{remarque}[Cas des corps finis]
Sur $\mathbb{F}_p$, le théorème de Fermat dit que $a^p = a$ pour tout $a \in \mathbb{F}_p$. Donc le polynôme $X^p - X$ est nul comme fonction, mais non nul comme polynôme (de degré $p$). Cette distinction est cruciale en théorie des corps finis et en codage.
\end{remarque}

\subsection{Polynômes à coefficients réels}

\begin{theoreme}[Racines complexes conjuguées]
Soit $P \in \mathbb{R}[X]$ et $z \in \mathbb{C}$ une racine de $P$ de multiplicité $m$. Alors $\bar{z}$ est aussi racine de $P$, avec la même multiplicité $m$.
\end{theoreme}

\begin{proof}
Comme $P \in \mathbb{R}[X]$, $\bar{P} = P$. Donc $P(\bar{z}) = \overline{P(z)} = \bar{0} = 0$. Pour la multiplicité : si $P = (X - z)^m Q$, alors en conjuguant $P = (X - \bar{z})^m \bar{Q}$, et la multiplicité de $\bar{z}$ est exactement $m$.
\end{proof}

\begin{corollaire}[Racine réelle d'un polynôme de degré impair]
Tout polynôme de $\mathbb{R}[X]$ de degré impair admet au moins une racine réelle.
\end{corollaire}

\begin{proof}
Sur $\mathbb{C}$, $P$ admet $\deg P$ racines comptées avec multiplicité (théorème fondamental, §1). Les racines non réelles s'apparient en couples $(z, \bar{z})$ de même multiplicité, donc leur nombre total est pair. Comme $\deg P$ est impair, le nombre de racines réelles l'est aussi : il y en a au moins une.

(Argument analytique alternatif : $\lim_{x \to \pm\infty} P(x) = \pm \infty$ avec des signes opposés ; le théorème des valeurs intermédiaires donne une racine.)
\end{proof}

\subsection{Exemples résolus (Difficulté ascendante)}

\textbf{Exemple 1 (Factorisation par racines évidentes --- Facile).}
Factoriser $P = X^3 - 6X^2 + 11X - 6$ dans $\mathbb{R}[X]$.

\textit{Correction :} on teste les diviseurs entiers de $6$. $P(1) = 1 - 6 + 11 - 6 = 0$, donc $1$ est racine. Division : $P = (X - 1)(X^2 - 5X + 6) = (X - 1)(X - 2)(X - 3)$.

\medskip

\textbf{Exemple 2 (Multiplicité par dérivées --- Intermédiaire).}
Soit $P = X^4 - 4X^3 + 6X^2 - 4X + 1$. Montrer que $1$ est racine de multiplicité $4$.

\textit{Correction :} on remarque que $P = (X - 1)^4$ par le binôme. Vérification par dérivées :
$P(1) = 0$. $P' = 4X^3 - 12X^2 + 12X - 4$, $P'(1) = 0$. $P'' = 12X^2 - 24X + 12$, $P''(1) = 0$. $P''' = 24X - 24$, $P'''(1) = 0$. $P^{(4)} = 24 \neq 0$. Donc $\mathrm{mult}_1(P) = 4$. \checkmark

\medskip

\textbf{Exemple 3 (Polynôme à coefficients réels et racine complexe --- Difficile).}
Le polynôme $P = X^4 + 1$ a-t-il des racines réelles ? Le factoriser dans $\mathbb{R}[X]$.

\textit{Correction :} pour tout $x \in \mathbb{R}$, $x^4 + 1 \ge 1 > 0$, donc $P$ n'a pas de racine réelle.
Sur $\mathbb{C}$, les racines sont les solutions de $z^4 = -1 = e^{i\pi}$, soit $z_k = e^{i(\pi + 2k\pi)/4}$ pour $k = 0, 1, 2, 3$ :
\[ z_0 = e^{i\pi/4},\ z_1 = e^{3i\pi/4},\ z_2 = e^{5i\pi/4} = \bar{z_1},\ z_3 = e^{7i\pi/4} = \bar{z_0}. \]
Pour factoriser sur $\mathbb{R}$, on regroupe les racines conjuguées :
\begin{align*}
(X - z_0)(X - \bar{z_0}) &= X^2 - 2 \mathrm{Re}(z_0) X + |z_0|^2 = X^2 - \sqrt{2} X + 1, \\
(X - z_1)(X - \bar{z_1}) &= X^2 - 2 \mathrm{Re}(z_1) X + 1 = X^2 + \sqrt{2} X + 1.
\end{align*}
Donc $P = (X^2 - \sqrt{2} X + 1)(X^2 + \sqrt{2} X + 1)$.

\subsection{Exercices pratiques avec Python}
\texttt{sympy} permet le calcul exact des racines avec multiplicité, la factorisation symbolique et la dérivation.

\begin{lstlisting}
import sympy as sp

X = sp.symbols('X')

# 1. Factorisation et racines avec multiplicite
P = X**5 - 5*X**4 + 10*X**3 - 10*X**2 + 5*X - 1
print(f"P = {P}")
print(f"  factorise   : {sp.factor(P)}")
print(f"  racines     : {sp.roots(P, X)}")  # dictionnaire racine -> multiplicite

# 2. Caracterisation par derivees : 1 est-il racine multiple ?
def multiplicite(poly, a, var):
    """Multiplicite de a dans poly, par derivees successives."""
    m = 0
    Q = sp.expand(poly)
    while sp.simplify(Q.subs(var, a)) == 0:
        m += 1
        Q = sp.diff(Q, var)
    return m

P_test = (X - 1)**3 * (X + 2)**2 * (X - 5)
print(f"\nP_test = {sp.expand(P_test)}")
for a in [-2, 1, 5, 0]:
    print(f"  mult({a}) = {multiplicite(P_test, a, X)}")

# 3. Coefficients reels => racines complexes conjuguees
P_real = X**4 + 1
print(f"\nRacines complexes de X^4 + 1 :")
for r in sp.solve(P_real, X):
    print(f"  {r}  (numerique : {complex(r):.4f})")
print(f"Factorisation reelle : {sp.factor(P_real, extension=[sp.sqrt(2)])}")
\end{lstlisting}

\subsection{Exercices à faire}
\begin{enumerate}
    \item \textbf{Exercice 1.} Factoriser $X^4 - 1$ dans $\mathbb{R}[X]$ et dans $\mathbb{C}[X]$. Combien de racines distinctes dans chaque corps ?
    \item \textbf{Exercice 2.} Trouver le polynôme unitaire de $\mathbb{R}[X]$ de degré $4$ admettant $1$ pour racine simple, $2$ pour racine double et $-1$ pour racine simple. Vérifier par développement.
    \item \textbf{Exercice 3.} Démontrer que tout polynôme de $\mathbb{R}[X]$ se factorise comme produit (à un scalaire près) de polynômes de degré $1$ et de polynômes de degré $2$ à discriminant strictement négatif. (Utiliser la factorisation sur $\mathbb{C}$ et le théorème des racines conjuguées.)
    \item \textbf{Exercice 4.} Soit $P \in \mathbb{R}[X]$ avec $\deg P = n \ge 1$. Montrer qu'entre deux racines réelles consécutives de $P$, $P'$ admet au moins une racine. (Théorème de Rolle algébrique --- on pourra invoquer Rolle de l'analyse.)
    \item \textbf{Exercice 5 (synthèse).} Sur le corps $\mathbb{F}_5$, lister tous les polynômes unitaires de degré $2$ qui n'ont aucune racine dans $\mathbb{F}_5$. (Indication : il y en a $10$.)
\end{enumerate}

\section{Arithmétique dans $K[X]$}

\subsection{Motivation}
L'anneau $K[X]$ partage avec l'anneau $\mathbb{Z}$ une propriété cruciale : il est muni d'une \emph{division euclidienne}. Cette analogie n'est pas anecdotique --- elle entraîne mécaniquement toute la théorie arithmétique du chapitre III (PGCD, algorithme d'Euclide, identité de Bézout, théorème de Gauss). En pratique, cela permet de \emph{factoriser} les polynômes, de \emph{simplifier les fractions rationnelles} (chapitre VI) et fonde l'arithmétique des codes correcteurs et des extensions de corps. Sur le plan algorithmique, le \emph{procédé de Hörner} fournit une évaluation et une division par $(X - a)$ en temps linéaire, technique anciennement attribuée à Sharaf Eddine Al-Tusi (XII$^{\text{e}}$ siècle).

\subsection{Approche par problème : PGCD de deux polynômes}
\textbf{Problème :} Calculer le PGCD de $A = X^4 - 1$ et $B = X^3 - X$ dans $\mathbb{R}[X]$.

\textbf{Résolution :} On peut deviner par factorisation : $A = (X-1)(X+1)(X^2+1)$ et $B = X(X-1)(X+1)$. Les facteurs communs sont $(X-1)(X+1) = X^2 - 1$. Mais en l'absence de factorisation évidente, comment procéder algorithmiquement ? Réponse : par divisions euclidiennes successives, exactement comme dans $\mathbb{Z}$. C'est l'algorithme d'Euclide pour les polynômes, que cette section formalise.

\subsection{Divisibilité dans $K[X]$}

\begin{definition}[Divisibilité]
Soient $A, B \in K[X]$. On dit que $B$ \textbf{divise} $A$, et on note $B \mid A$, s'il existe $Q \in K[X]$ tel que $A = B Q$. Tout diviseur de $A$ est aussi appelé un \emph{facteur} de $A$.
\end{definition}

\begin{proposition}[Propriétés de la divisibilité]
Pour $A, B, C \in K[X]$ :
\begin{enumerate}
    \item $A \mid A$ (réflexivité) ; $1 \mid A$ ; $A \mid 0$.
    \item Si $A \mid B$ et $B \mid C$, alors $A \mid C$ (transitivité).
    \item Si $A \mid B$ et $A \mid C$, alors $A \mid (\lambda B + \mu C)$ pour tous $\lambda, \mu \in K[X]$.
    \item Si $A \mid B$ avec $A, B$ non nuls, alors $\deg A \le \deg B$.
\end{enumerate}
\end{proposition}

\begin{definition}[Polynômes associés]
Deux polynômes $A, B$ sont \textbf{associés} s'il existe $\lambda \in K^*$ tel que $A = \lambda B$. Cette relation est une équivalence ; deux polynômes associés ont les mêmes diviseurs.
\end{definition}

\begin{remarque}
Pour avoir l'unicité du PGCD (à venir), on impose souvent qu'il soit \emph{unitaire} (coefficient dominant $1$). C'est l'analogue dans $K[X]$ du choix « positif » dans $\mathbb{Z}$.
\end{remarque}

\subsection{Division euclidienne}

\begin{theoreme}[{Division euclidienne dans $K[X]$}]
Soient $A, B \in K[X]$ avec $B \neq 0$. Il existe un \emph{unique} couple $(Q, R) \in K[X]^2$ tel que :
\[ A = B Q + R \quad \text{avec} \quad \deg R < \deg B. \]
$Q$ est appelé le \textbf{quotient} et $R$ le \textbf{reste} de la division euclidienne de $A$ par $B$.
\end{theoreme}

\begin{proof}
\emph{Existence.} Récurrence sur $\deg A$. Si $\deg A < \deg B$, on prend $Q = 0$ et $R = A$. Sinon, posons $A = a_n X^n + \cdots$, $B = b_p X^p + \cdots$ avec $n \ge p$. Alors $A_1 = A - \frac{a_n}{b_p} X^{n-p} B$ a un degré strictement inférieur à $\deg A$. Par hypothèse de récurrence, $A_1 = B Q_1 + R$ avec $\deg R < \deg B$, d'où $A = B (Q_1 + \frac{a_n}{b_p} X^{n-p}) + R$.

\emph{Unicité.} Si $A = BQ_1 + R_1 = BQ_2 + R_2$ avec $\deg R_i < \deg B$, alors $B(Q_1 - Q_2) = R_2 - R_1$. Si $Q_1 \neq Q_2$, on aurait $\deg(B(Q_1 - Q_2)) \ge \deg B > \deg(R_2 - R_1)$, absurde. Donc $Q_1 = Q_2$ et $R_1 = R_2$.
\end{proof}

\begin{remarque}
Cette division se pratique « à la main » de manière analogue à la division des entiers en colonne : on aligne les puissances de $X$ et on élimine le terme dominant pas à pas.
\end{remarque}

\subsection{Procédé de Hörner et division par $(X - a)$}

\begin{proposition}[Procédé de Hörner]
Soit $P = a_n X^n + a_{n-1} X^{n-1} + \cdots + a_0 \in K[X]$ et $a \in K$. Posons :
\[ b_n = a_n, \quad b_{k-1} = a_{k-1} + a\, b_k \quad \text{pour } k = n, n-1, \dots, 1. \]
Alors :
\begin{itemize}
    \item $P(a) = b_0$ ;
    \item $Q(X) = b_n X^{n-1} + b_{n-1} X^{n-2} + \cdots + b_1$ est le quotient de la division euclidienne de $P$ par $(X - a)$, et $b_0$ est le reste : $P(X) = (X - a) Q(X) + b_0$.
\end{itemize}
Le coût est de $n$ multiplications et $n$ additions, contre $\sim n^2 / 2$ pour une évaluation naïve.
\end{proposition}

\begin{proof}
On procède par identification des coefficients dans $P(X) = (X - a)(b_n X^{n-1} + \cdots + b_1) + b_0$. Le coefficient de $X^k$ donne $a_k = b_k - a\, b_{k+1}$, soit $b_{k+1} = (a_k + b_k - \dots)$. La récurrence ascendante mène aux formules ci-dessus.
\end{proof}

\paragraph{Note historique : Sharaf Eddine Al-Tusi (vers 1135--1213).}
Mathématicien et astronome persan né à Tus (Iran), Sharaf al-D\={\i}n al-T\={u}s\={\i} a développé dans son traité sur les équations cubiques (\emph{Risāla fī al-mu‘ādalāt}) une méthode tabulaire d'évaluation et de division de polynômes équivalente au procédé que William G.\ Horner publiera en 1819. Sa technique est aussi utilisée pour calculer numériquement les racines d'une équation polynomiale par approximations successives. Cette même méthode était également connue des mathématiciens chinois (Qin Jiushao, XIII$^{\text{e}}$ siècle) et fut redécouverte par l'italien Paolo Ruffini en 1804. La dénomination plus juste serait donc \emph{procédé d'Al-Tusi--Ruffini--Horner}.

\subsection{PGCD dans $K[X]$}

\begin{definition}[Diviseur commun, PGCD]
Soient $A, B \in K[X]$ non tous deux nuls. Un \emph{diviseur commun} de $A$ et $B$ est un polynôme $D$ tel que $D \mid A$ et $D \mid B$.

Le \textbf{plus grand commun diviseur} (PGCD) de $A$ et $B$ est l'unique diviseur commun \emph{unitaire} de degré maximal. On le note $A \wedge B$ ou $\mathrm{pgcd}(A, B)$.
\end{definition}

\begin{theoreme}[Algorithme d'Euclide pour les polynômes]
Soient $A, B \in K[X]$ avec $B \neq 0$. La suite des restes obtenue par divisions euclidiennes successives :
\[ A = B Q_1 + R_1,\ B = R_1 Q_2 + R_2,\ R_1 = R_2 Q_3 + R_3,\ \dots \]
avec $\deg R_1 > \deg R_2 > \cdots$ se termine en un nombre fini d'étapes par un reste $R_N = 0$. Le dernier reste non nul $R_{N-1}$, normalisé pour être unitaire, est le PGCD de $A$ et $B$.
\end{theoreme}

\begin{proof}
La suite des degrés $\deg R_k$ est strictement décroissante, donc l'algorithme s'arrête. À chaque étape, l'égalité $R_{k-1} = R_k Q_{k+1} + R_{k+1}$ montre que $\{$diviseurs communs de $R_{k-1}, R_k\} = \{$diviseurs communs de $R_k, R_{k+1}\}$. Par récurrence, $\mathrm{pgcd}(A, B) = \mathrm{pgcd}(R_{N-1}, 0) = R_{N-1}$ (à un scalaire près).
\end{proof}

\subsection{Identité de Bézout}

\begin{theoreme}[{Bézout pour $K[X]$}]
Soient $A, B \in K[X]$ non tous deux nuls. Il existe $U, V \in K[X]$ tels que :
\[ A U + B V = A \wedge B. \]
\end{theoreme}

\begin{proof}
Considérons l'ensemble $I = \{ A U + B V : U, V \in K[X]\}$. C'est un \emph{idéal} de $K[X]$ contenant $A$ et $B$. Comme $K[X]$ est principal (chaque idéal est engendré par un polynôme), il existe $D \in K[X]$ tel que $I = (D)$. Ce $D$ divise $A$ et $B$, et tout autre diviseur commun divise $D$ : c'est le PGCD (à association près). De plus $D \in I$ s'écrit $D = A U + B V$.
\end{proof}

\begin{corollaire}[Polynômes premiers entre eux]
$A, B$ sont \textbf{premiers entre eux}, noté $A \wedge B = 1$, si et seulement si :
\[ \exists U, V \in K[X],\quad A U + B V = 1. \]
\end{corollaire}

\subsection{Théorème de Gauss}

\begin{theoreme}[Gauss]
Soient $A, B, C \in K[X]$. Si $A \mid B C$ et $A \wedge B = 1$, alors $A \mid C$.
\end{theoreme}

\begin{proof}
Bézout : $A U + B V = 1$, donc $C = A U C + B V C$. Comme $A \mid A U C$ et $A \mid B C V$ (puisque $A \mid B C$), on a $A \mid C$.
\end{proof}

\begin{corollaire}
\begin{enumerate}
    \item Si $A \wedge B = 1$ et $A \mid C$ et $B \mid C$, alors $A B \mid C$.
    \item Si $A \wedge B_1 = 1$ et $A \wedge B_2 = 1$, alors $A \wedge (B_1 B_2) = 1$.
\end{enumerate}
\end{corollaire}

\subsection{PPCM}

\begin{definition}[PPCM]
Le \textbf{plus petit commun multiple} de $A$ et $B$ non nuls, noté $A \vee B$, est l'unique polynôme unitaire de degré minimal tel que $A \mid M$ et $B \mid M$.
\end{definition}

\begin{proposition}[Lien PGCD/PPCM]
Pour $A, B$ unitaires non nuls : $A \cdot B = (A \wedge B) \cdot (A \vee B)$.
\end{proposition}

\subsection{Exemples résolus (Difficulté ascendante)}

\textbf{Exemple 1 (Division euclidienne directe --- Facile).}
Effectuer la division euclidienne de $A = X^4 + 2X^3 - X + 1$ par $B = X^2 + X - 1$ dans $\mathbb{R}[X]$.

\textit{Correction :} on procède en colonnes.
$X^4 \div X^2 = X^2$ ; $X^2 \cdot B = X^4 + X^3 - X^2$. Reste partiel : $A - X^2 B = X^3 + X^2 - X + 1$.
$X^3 \div X^2 = X$ ; $X \cdot B = X^3 + X^2 - X$. Reste : $1$.
Le quotient est $Q = X^2 + X$ et le reste est $R = 1$. Vérification : $(X^2 + X)(X^2 + X - 1) + 1 = X^4 + 2X^3 - X + 1$. \checkmark

\medskip

\textbf{Exemple 2 (Hörner --- Intermédiaire).}
Évaluer $P = 3X^4 - 2X^3 + 5X - 7$ en $a = 2$, et donner le quotient de $P$ par $(X - 2)$.

\textit{Correction :} on tabule les coefficients $a_4 = 3,\ a_3 = -2,\ a_2 = 0,\ a_1 = 5,\ a_0 = -7$.
\begin{center}
\begin{tabular}{c|cccccc}
 & $a_4$ & $a_3$ & $a_2$ & $a_1$ & $a_0$ \\
 & 3 & $-2$ & 0 & 5 & $-7$ \\
\hline
$\times 2$ : & & 6 & 8 & 16 & 42 \\
\hline
$b_k$ : & 3 & 4 & 8 & 21 & 35
\end{tabular}
\end{center}
Donc $P(2) = 35$ et $Q(X) = 3X^3 + 4X^2 + 8X + 21$, c'est-à-dire $P(X) = (X - 2)(3X^3 + 4X^2 + 8X + 21) + 35$.

\medskip

\textbf{Exemple 3 (Algorithme d'Euclide étendu --- Difficile).}
Calculer $A \wedge B$ pour $A = X^4 + X^3 - X - 1$ et $B = X^3 - 1$ dans $\mathbb{R}[X]$, puis trouver une identité de Bézout.

\textit{Correction :} divisions successives.
\[ A = X^4 + X^3 - X - 1 = (X + 1)(X^3 - 1) + 0 \cdot X^2 + 0 \cdot X + 0 \quad ?\ \text{Vérifions} : (X+1)(X^3-1) = X^4 - X + X^3 - 1 = X^4 + X^3 - X - 1.\ \checkmark \]
Donc $A = (X+1) B + 0$, le reste est nul, et $A \wedge B = $ partie unitaire de $B = X^3 - 1$.

Identité de Bézout : $A \cdot 1 + B \cdot (-(X+1)) = 0$ ne fonctionne pas. Comme $B \mid A$, le PGCD est $B$ (unitaire) et l'on peut écrire $B = A \cdot 0 + B \cdot 1$, soit $U = 0, V = 1$.

\medskip

\noindent\textbf{Variante.} Pour $A = X^3 + 1$ et $B = X^2 + X$ :
\[ A = B \cdot X + (-X^2 + 1) \quad \text{(reste } R_1 = -X^2 + 1\text{)} \]
\[ B = R_1 \cdot (-1) + (X + 1) \quad \text{(reste } R_2 = X + 1\text{)} \]
\[ R_1 = R_2 \cdot (-X + 1) + 0 \]
Donc $A \wedge B = X + 1$ (déjà unitaire). En remontant les égalités :
$X + 1 = B + R_1 = B + (A - X B) = A - (X - 1) B$. Soit $U = 1$ et $V = -(X - 1) = 1 - X$ : $A \cdot 1 + B \cdot (1 - X) = X + 1$.

\subsection{Exercices pratiques avec Python}
\texttt{sympy} fournit \texttt{div}, \texttt{gcd}, \texttt{gcdex} (Bézout étendu) sur les polynômes.

\begin{lstlisting}
import sympy as sp

X = sp.symbols('X')

# 1. Division euclidienne
A = sp.Poly(X**4 + 2*X**3 - X + 1, X)
B = sp.Poly(X**2 + X - 1, X)
Q, R = sp.div(A, B, X)
print(f"A / B :  Q = {Q.as_expr()},  R = {R.as_expr()}")

# 2. Procede de Horner (evaluation + division par (X - a))
def horner(coeffs, a):
    """coeffs : ordre decroissant des coefficients. Retourne (P(a), coeffs du quotient)."""
    bs = []
    b = 0
    for c in coeffs:
        b = b * a + c
        bs.append(b)
    return bs[-1], bs[:-1]    # (P(a), coefficients du quotient en ordre decroissant)

Pa, coeffs_Q = horner([3, -2, 0, 5, -7], 2)
print(f"\nHorner sur P=3X^4-2X^3+5X-7 en a=2 :")
print(f"  P(2) = {Pa}")
print(f"  Q(X) coefficients (ordre decroissant) = {coeffs_Q}")

# 3. Algorithme d'Euclide etendu : Bezout
A = sp.Poly(X**3 + 1, X)
B = sp.Poly(X**2 + X, X)
g = sp.gcd(A, B)
U, V, d = sp.gcdex(A.as_expr(), B.as_expr(), X)
print(f"\nA = {A.as_expr()},  B = {B.as_expr()}")
print(f"  PGCD     = {g.as_expr()}")
print(f"  Bezout   : ({U}) * A + ({V}) * B = {sp.expand(U * A.as_expr() + V * B.as_expr())}")

# 4. PPCM
print(f"\nPPCM(A, B) = {sp.lcm(A, B).as_expr()}")
print(f"Verification A*B / pgcd = {sp.expand(A.as_expr() * B.as_expr() / g.as_expr())}")
\end{lstlisting}

\subsection{Exercices à faire}
\begin{enumerate}
    \item \textbf{Exercice 1.} Effectuer la division euclidienne de $A = X^5 - X + 1$ par $B = X^2 + 1$ dans $\mathbb{R}[X]$.
    \item \textbf{Exercice 2.} Évaluer $P = 2X^4 - 3X^3 + X^2 - 5X + 6$ en $a = -1$ par le procédé de Hörner. En déduire le quotient de $P$ par $(X + 1)$ et $P(-1)$.
    \item \textbf{Exercice 3.} Calculer $A \wedge B$ pour $A = X^5 + X^4 + X^3 + X^2 + X + 1$ et $B = X^4 - 1$. Donner une identité de Bézout.
    \item \textbf{Exercice 4.} Démontrer que pour tous $m, n \in \mathbb{N}^*$, $\mathrm{pgcd}(X^m - 1, X^n - 1) = X^{\mathrm{pgcd}(m, n)} - 1$ dans $K[X]$. (Indication : utiliser la division euclidienne sur les exposants, à l'image du calcul du PGCD entier.)
    \item \textbf{Exercice 5 (synthèse).} Soit $A, B, C \in K[X]$. Montrer que si $A \wedge B = 1$ et $A \mid C^k$ pour un certain $k \ge 1$, et si $A$ est irréductible, alors $A \mid C$. (Variante du théorème de Gauss --- démontrée en §4 dans le cas général.)
\end{enumerate}

\section{Polynômes irréductibles, décomposition, relations racines--coefficients}

\subsection{Motivation}
Dans $\mathbb{Z}$, les nombres premiers sont les briques élémentaires : tout entier se décompose de manière unique en facteurs premiers. Dans $K[X]$, ce rôle est joué par les \textbf{polynômes irréductibles}, et la décomposition unique en irréductibles permet de comprendre la structure complète d'un polynôme : ses racines, ses facteurs réels, ses simplifications. Les polynômes irréductibles sur les corps finis $\mathbb{F}_p$ servent à construire les corps $\mathbb{F}_{p^n}$, fondement des codes correcteurs et de la cryptographie sur courbes elliptiques. Pour finir, les \emph{relations racines--coefficients} (formules de Viète) donnent un dictionnaire entre les coefficients d'un polynôme et les fonctions symétriques de ses racines.

\subsection{Approche par problème : où $X^2 + 1$ est-il irréductible ?}
\textbf{Problème :} Le polynôme $P = X^2 + 1$ est-il irréductible dans $\mathbb{Q}[X]$ ? dans $\mathbb{R}[X]$ ? dans $\mathbb{C}[X]$ ?

\textbf{Résolution :}
\begin{itemize}
    \item Sur $\mathbb{C}$ : $X^2 + 1 = (X - i)(X + i)$. \emph{Réductible.}
    \item Sur $\mathbb{R}$ : aucune racine réelle (puisque $x^2 + 1 \ge 1$). Si $P = AB$ avec $A, B \in \mathbb{R}[X]$ non constants, alors $\deg A = \deg B = 1$, donc $P$ aurait deux racines réelles --- contradiction. \emph{Irréductible.}
    \item Sur $\mathbb{Q}$ : aucune racine rationnelle (a fortiori, pas de racine réelle). Le même argument donne l'irréductibilité. \emph{Irréductible.}
\end{itemize}

L'irréductibilité \emph{dépend du corps de base}. Cette section caractérise les irréductibles dans $\mathbb{C}[X]$ et $\mathbb{R}[X]$.

\subsection{Polynômes irréductibles}

\begin{definition}[Polynôme irréductible]
Un polynôme $P \in K[X]$ est \textbf{irréductible} sur $K$ si :
\begin{enumerate}
    \item $\deg P \ge 1$ (on exclut les constantes) ;
    \item dans toute factorisation $P = A \cdot B$, l'un des facteurs $A$ ou $B$ est inversible dans $K[X]$, c'est-à-dire constant non nul.
\end{enumerate}
Un polynôme non constant qui n'est pas irréductible est dit \textbf{réductible}.
\end{definition}

\begin{exemple}
\begin{itemize}
    \item Tout polynôme de degré $1$ est irréductible (sur n'importe quel corps).
    \item $X^2 - 2$ est irréductible sur $\mathbb{Q}$ mais réductible sur $\mathbb{R}$ ($X^2 - 2 = (X - \sqrt{2})(X + \sqrt{2})$).
    \item $X^2 + 1$ est irréductible sur $\mathbb{R}$ mais réductible sur $\mathbb{C}$.
    \item Sur $\mathbb{F}_2$, le polynôme $X^2 + X + 1$ est irréductible (pas de racine dans $\mathbb{F}_2$). Il sert à construire $\mathbb{F}_4$.
\end{itemize}
\end{exemple}

\subsection{Décomposition unique en facteurs irréductibles}

\begin{theoreme}[{Théorème fondamental de l'arithmétique dans $K[X]$}]
Tout polynôme non constant $P \in K[X]$ se factorise sous la forme
\[ P = c \cdot P_1^{m_1} \cdot P_2^{m_2} \cdots P_r^{m_r}, \]
où $c \in K^*$ et les $P_i$ sont des polynômes irréductibles \emph{unitaires} deux à deux distincts. Cette décomposition est \emph{unique} à l'ordre près des facteurs.
\end{theoreme}

\begin{proof}[Esquisse]
\emph{Existence} : par récurrence sur $\deg P$. Si $P$ est irréductible, c'est terminé. Sinon $P = AB$ avec $\deg A, \deg B < \deg P$, et chaque facteur se décompose par récurrence.

\emph{Unicité} : utilise le théorème de Gauss. Si $P_1 \cdots P_r = Q_1 \cdots Q_s$ avec tous les facteurs irréductibles unitaires, alors $P_1$ divise $Q_1 \cdots Q_s$. Comme $P_1$ est irréductible et premier avec tout $Q_j$ qui ne lui est pas associé, Gauss force $P_1$ à être associé à un certain $Q_j$. On simplifie et l'on conclut par récurrence.
\end{proof}

\subsection{Irréductibles de $\mathbb{C}[X]$}

\begin{theoreme}[Théorème fondamental de l'algèbre, admis]
Tout polynôme non constant de $\mathbb{C}[X]$ admet au moins une racine dans $\mathbb{C}$. Autrement dit, $\mathbb{C}$ est \textbf{algébriquement clos}.
\end{theoreme}

\begin{remarque}
Ce théorème, conjecturé par d'Alembert et démontré rigoureusement par Gauss en 1799 (sa thèse), nécessite des outils d'analyse. Il sera démontré en cours de fonctions complexes.
\end{remarque}

\begin{corollaire}[{Irréductibles de $\mathbb{C}[X]$}]
Les polynômes irréductibles de $\mathbb{C}[X]$ sont exactement ceux de degré $1$.
\end{corollaire}

\begin{corollaire}[{Décomposition dans $\mathbb{C}[X]$}]
Tout $P \in \mathbb{C}[X]$ de degré $n \ge 1$ et de coefficient dominant $a_n$ se factorise de manière unique :
\[ P = a_n \prod_{k=1}^{r} (X - z_k)^{m_k}, \qquad \sum_{k=1}^{r} m_k = n, \]
où les $z_k \in \mathbb{C}$ sont les racines distinctes de $P$ et $m_k$ leurs multiplicités.
\end{corollaire}

\subsection{Irréductibles de $\mathbb{R}[X]$}

\begin{theoreme}[{Irréductibles de $\mathbb{R}[X]$}]
Les polynômes irréductibles de $\mathbb{R}[X]$ sont :
\begin{itemize}
    \item les polynômes de degré $1$ ;
    \item les polynômes de degré $2$ à \emph{discriminant strictement négatif}, c'est-à-dire de la forme $aX^2 + bX + c$ avec $a \neq 0$ et $b^2 - 4ac < 0$.
\end{itemize}
\end{theoreme}

\begin{proof}
Les polynômes de degré $1$ sont irréductibles. Pour le degré $2$ : un trinôme $aX^2 + bX + c$ a une racine réelle ssi $b^2 - 4ac \ge 0$. S'il en a une, il se factorise en $a(X - r_1)(X - r_2)$, donc réductible. Sinon, ses deux racines complexes sont conjuguées et il est irréductible sur $\mathbb{R}$.

Réciproquement, soit $P \in \mathbb{R}[X]$ de degré $\ge 3$. Sur $\mathbb{C}$, il admet une racine $z$. Si $z \in \mathbb{R}$, alors $(X - z) \mid P$ dans $\mathbb{R}[X]$, donc $P$ est réductible. Sinon, $\bar{z}$ est aussi racine (conjugaison), et $(X - z)(X - \bar{z}) = X^2 - 2 \mathrm{Re}(z) X + |z|^2 \in \mathbb{R}[X]$ divise $P$ dans $\mathbb{R}[X]$ ; encore réductible.
\end{proof}

\begin{corollaire}[{Décomposition dans $\mathbb{R}[X]$}]
Tout $P \in \mathbb{R}[X]$ de degré $n \ge 1$ se factorise de manière unique sous la forme
\[ P = a_n \prod_{i=1}^{p} (X - a_i)^{\alpha_i} \prod_{j=1}^{q} (X^2 + b_j X + c_j)^{\beta_j}, \]
où les $a_i \in \mathbb{R}$ sont les racines réelles, $\alpha_i$ leurs multiplicités, et les facteurs $X^2 + b_j X + c_j$ sont irréductibles ($b_j^2 - 4 c_j < 0$).
\end{corollaire}

\subsection{Relations entre coefficients et racines (formules de Viète)}

Soient $z_1, \dots, z_n$ des indéterminées. Les \textbf{polynômes symétriques élémentaires} sont définis par :
\begin{align*}
\sigma_1 &= z_1 + z_2 + \cdots + z_n, \\
\sigma_2 &= \sum_{i < j} z_i z_j, \\
\sigma_k &= \sum_{i_1 < i_2 < \dots < i_k} z_{i_1} z_{i_2} \cdots z_{i_k}, \\
\sigma_n &= z_1 z_2 \cdots z_n.
\end{align*}

\begin{theoreme}[Formules de Viète]
Soit $P = a_n X^n + a_{n-1} X^{n-1} + \cdots + a_0 \in K[X]$ un polynôme \emph{scindé} sur $K$, c'est-à-dire admettant $n$ racines $z_1, \dots, z_n \in K$ (comptées avec multiplicité). Alors :
\[ \boxed{\ \sigma_k(z_1, \dots, z_n) = (-1)^k\, \frac{a_{n-k}}{a_n}\ } \quad \text{pour } k = 1, 2, \dots, n. \]
\end{theoreme}

\begin{proof}
On part de $P = a_n (X - z_1)(X - z_2) \cdots (X - z_n)$ et l'on développe. Le coefficient de $X^{n-k}$ dans le produit est $(-1)^k \sigma_k$ multiplié par $a_n$, soit $a_{n-k} = a_n \cdot (-1)^k \sigma_k$.
\end{proof}

\begin{remarque}[Cas usuels]
Pour le degré 2 ($P = aX^2 + bX + c$ de racines $z_1, z_2$) :
\[ z_1 + z_2 = -\frac{b}{a}, \qquad z_1 z_2 = \frac{c}{a}. \]
Pour le degré 3 ($P = X^3 + pX^2 + qX + r$ unitaire de racines $z_1, z_2, z_3$) :
\[ z_1 + z_2 + z_3 = -p, \quad z_1 z_2 + z_1 z_3 + z_2 z_3 = q, \quad z_1 z_2 z_3 = -r. \]
\end{remarque}

\begin{remarque}
Sur $\mathbb{C}$, tout polynôme est scindé (par le théorème fondamental). Les formules de Viète s'appliquent donc toujours dans $\mathbb{C}[X]$.
\end{remarque}

\subsection{Exemples résolus (Difficulté ascendante)}

\textbf{Exemple 1 (Décomposition dans $\mathbb{R}$ et $\mathbb{C}$ --- Facile).}
Décomposer $P = X^4 - 1$ dans $\mathbb{R}[X]$ et dans $\mathbb{C}[X]$.

\textit{Correction :}
$P = (X^2 - 1)(X^2 + 1) = (X - 1)(X + 1)(X^2 + 1)$ dans $\mathbb{R}[X]$ : $X^2 + 1$ est irréductible sur $\mathbb{R}$ (discriminant $-4 < 0$).
Sur $\mathbb{C}$ : $X^2 + 1 = (X - i)(X + i)$, d'où $P = (X - 1)(X + 1)(X - i)(X + i)$. Quatre facteurs de degré $1$.

\medskip

\textbf{Exemple 2 (Application de Viète --- Intermédiaire).}
Soient $z_1, z_2, z_3$ les racines complexes de $P = X^3 - 6X^2 + 11X - 6$. Calculer $z_1 + z_2 + z_3$, $z_1 z_2 z_3$, et $z_1^2 + z_2^2 + z_3^2$.

\textit{Correction :} Viète donne directement $\sigma_1 = z_1 + z_2 + z_3 = 6$ et $\sigma_3 = z_1 z_2 z_3 = 6$, ainsi que $\sigma_2 = z_1 z_2 + z_1 z_3 + z_2 z_3 = 11$.
Pour la somme des carrés, on utilise $z_1^2 + z_2^2 + z_3^2 = \sigma_1^2 - 2 \sigma_2 = 36 - 22 = 14$.

(Vérification : on peut factoriser $P = (X-1)(X-2)(X-3)$, donc les racines sont $1, 2, 3$ avec $1 + 4 + 9 = 14$. \checkmark)

\medskip

\textbf{Exemple 3 (Décomposition réelle complète --- Difficile).}
Décomposer $P = X^6 - 1$ dans $\mathbb{R}[X]$ et dans $\mathbb{C}[X]$.

\textit{Correction :} sur $\mathbb{C}$, les racines sont les racines $6$-ièmes de l'unité : $\omega^k = e^{2 i k \pi / 6}$ pour $k = 0, 1, \dots, 5$. Donc
\[ X^6 - 1 = \prod_{k=0}^{5} (X - e^{i k \pi / 3}). \]
On les liste : $1,\ e^{i\pi/3},\ e^{2i\pi/3},\ -1,\ e^{4i\pi/3},\ e^{5i\pi/3}$. Les conjuguées s'apparient : $(e^{i\pi/3}, e^{5i\pi/3})$ et $(e^{2i\pi/3}, e^{4i\pi/3})$.
Sur $\mathbb{R}$, on regroupe les conjuguées :
\begin{align*}
(X - e^{i\pi/3})(X - e^{-i\pi/3}) &= X^2 - 2\cos(\pi/3) X + 1 = X^2 - X + 1, \\
(X - e^{2i\pi/3})(X - e^{-2i\pi/3}) &= X^2 - 2\cos(2\pi/3) X + 1 = X^2 + X + 1.
\end{align*}
D'où :
\[ X^6 - 1 = (X - 1)(X + 1)(X^2 - X + 1)(X^2 + X + 1). \]
On vérifie que les deux trinômes ont des discriminants $1 - 4 = -3 < 0$ : ils sont bien irréductibles sur $\mathbb{R}$.

\subsection{Exercices pratiques avec Python}
\texttt{sympy.factor} effectue la factorisation symbolique sur $\mathbb{Q}$, et accepte des extensions algébriques pour factoriser sur $\mathbb{R}$ ou $\mathbb{C}$. Les formules de Viète se vérifient via \texttt{sympy.elementary\_symmetric\_polynomial}.

\begin{lstlisting}
import sympy as sp
from sympy import I, sqrt, factor, roots, expand

X = sp.symbols('X')

# 1. Decomposition selon le corps de base
P = X**6 - 1
print(f"X^6 - 1 sur Q :  {factor(P)}")
print(f"  sur Q(sqrt(3)) : {factor(P, extension=[sqrt(3)])}")
print(f"  sur C          : {factor(P, extension=[I, sqrt(3)])}")

# 2. Verification des formules de Viete
P = X**3 - 6*X**2 + 11*X - 6
zs = list(roots(P).keys())
print(f"\nRacines de {P} : {zs}")
print(f"  sigma_1 = {sum(zs)}        (attendu : 6)")
print(f"  sigma_2 = {sum(z1*z2 for i, z1 in enumerate(zs) for z2 in zs[i+1:])}   (attendu : 11)")
print(f"  sigma_3 = {sp.prod(zs)}    (attendu : 6)")

# 3. Somme des carres a partir de sigma_1, sigma_2 (sans calculer les racines)
P = X**3 - 6*X**2 + 11*X - 6
coeffs = sp.Poly(P, X).all_coeffs()
a3, a2, a1, a0 = coeffs
sigma1 = -a2 / a3
sigma2 =  a1 / a3
somme_carres = sigma1**2 - 2 * sigma2
print(f"\nz1^2 + z2^2 + z3^2 = sigma_1^2 - 2 sigma_2 = {somme_carres}")

# 4. Test d'irreductibilite sur Q (sans factoriser explicitement)
for poly in [X**2 + 1, X**2 - 2, X**3 - 2, X**4 + 1, X**2 + X + 1]:
    print(f"{poly} irreductible sur Q ? {sp.Poly(poly, X).is_irreducible}")
\end{lstlisting}

\subsection{Exercices à faire}
\begin{enumerate}
    \item \textbf{Exercice 1.} Décomposer $P = X^4 + 1$ dans $\mathbb{Q}[X]$, $\mathbb{R}[X]$ et $\mathbb{C}[X]$. Vérifier que sur $\mathbb{Q}$, $P$ est irréductible (par exemple par le critère d'Eisenstein appliqué à $P(X+1)$, ou par considération de degré).
    \item \textbf{Exercice 2.} Soient $a, b, c$ les racines complexes de $P = X^3 + 2X^2 - X + 5$. Calculer $a + b + c$, $ab + ac + bc$, $abc$, puis $a^2 + b^2 + c^2$ et $\frac{1}{a} + \frac{1}{b} + \frac{1}{c}$ (cette dernière sans expliciter les racines).
    \item \textbf{Exercice 3.} Trouver un polynôme unitaire $P \in \mathbb{R}[X]$ de degré $3$ dont les racines sont $2$, $1 + i$ et $1 - i$. Vérifier les formules de Viète.
    \item \textbf{Exercice 4.} Décomposer $P = X^4 + X^2 + 1$ dans $\mathbb{R}[X]$. (Indication : remarquer que $P = (X^2 + X + 1)(X^2 - X + 1)$ par identité algébrique ; vérifier en développant.)
    \item \textbf{Exercice 5 (synthèse).} Soit $P \in \mathbb{R}[X]$ unitaire de degré $n$, scindé sur $\mathbb{R}$ avec racines réelles $a_1 \le a_2 \le \cdots \le a_n$. Démontrer que $\frac{P'}{P} = \sum_{i=1}^{n} \frac{1}{X - a_i}$ (en tant qu'identité de fractions rationnelles --- voir chapitre VI).
\end{enumerate}

\section*{Travaux Pratiques (TP) Python}
\addcontentsline{toc}{section}{Travaux Pratiques (TP) Python}

\vspace{0.3cm}
\noindent L'objectif de ce TP est de construire pas à pas une bibliothèque \emph{maison} de manipulation de polynômes : représentation, arithmétique, racines, factorisation. Ce travail prépare à la décomposition en éléments simples (chapitre VI) et à des applications concrètes en codes correcteurs et en cryptographie. On comparera systématiquement à \texttt{numpy.polynomial} et \texttt{sympy}.

\vspace{0.3cm}

\noindent\textbf{Exercice 1 : Classe \texttt{Polynome}}
\begin{enumerate}
    \item \textbf{[Bloom : Compréhension | Difficulté : Moyenne]} \\
    Créer une classe \texttt{Polynome} dont l'attribut principal est la liste des coefficients en ordre croissant (\texttt{[a0, a1, ..., an]}). Implémenter :
    \begin{itemize}
        \item \texttt{\_\_init\_\_(self, coeffs)} avec normalisation (suppression des zéros de tête) ;
        \item \texttt{degre(self)} ;
        \item \texttt{\_\_repr\_\_(self)} pour afficher joliment le polynôme.
    \end{itemize}
    \item \textbf{[Bloom : Application | Difficulté : Moyenne]} \\
    Surcharger les opérateurs \texttt{\_\_add\_\_}, \texttt{\_\_sub\_\_}, \texttt{\_\_mul\_\_} pour la somme, la différence et le produit (produit de Cauchy).
    \item \textbf{[Bloom : Évaluation | Difficulté : Moyenne]} \\
    Tester votre classe sur $P = 2X^3 - X + 5$ et $Q = X^2 + 3$ : vérifier $\deg(PQ) = \deg P + \deg Q$.
\end{enumerate}

\vspace{0.3cm}

\noindent\textbf{Exercice 2 : Évaluation par le procédé de Hörner}
\begin{enumerate}
    \item \textbf{[Bloom : Application | Difficulté : Moyenne]} \\
    Ajouter une méthode \texttt{evaluer(self, a)} qui calcule $P(a)$ par le procédé de Hörner (en $n$ multiplications et $n$ additions).
    \item \textbf{[Bloom : Évaluation | Difficulté : Moyenne]} \\
    Comparer le temps d'exécution de \texttt{evaluer\_Horner} et d'une évaluation naïve $\sum a_k a^k$ (avec puissances explicites) pour un polynôme de degré $1000$ évalué $10000$ fois. Conclusion ?
    \item \textbf{[Bloom : Synthèse | Difficulté : Moyenne]} \\
    Modifier le procédé pour qu'il retourne en plus le quotient $Q$ tel que $P = (X - a) Q + P(a)$ (division par $(X - a)$ par Hörner).
\end{enumerate}

\vspace{0.3cm}

\noindent\textbf{Exercice 3 : Division euclidienne dans $K[X]$}
\begin{enumerate}
    \item \textbf{[Bloom : Synthèse | Difficulté : Difficile]} \\
    Implémenter \texttt{divmod\_poly(A, B)} qui retourne le couple $(Q, R)$ de la division euclidienne $A = BQ + R$ avec $\deg R < \deg B$. Suivre la démonstration du cours : éliminer pas à pas le terme dominant.
    \item \textbf{[Bloom : Évaluation | Difficulté : Moyenne]} \\
    Vérifier sur $A = X^4 + 2X^3 - X + 1$ et $B = X^2 + X - 1$ que $Q = X^2 + X$ et $R = 1$.
    \item \textbf{[Bloom : Évaluation | Difficulté : Moyenne]} \\
    Comparer avec \texttt{sympy.div(A, B, X)}.
\end{enumerate}

\vspace{0.3cm}

\noindent\textbf{Exercice 4 : PGCD de polynômes (algorithme d'Euclide)}
\begin{enumerate}
    \item \textbf{[Bloom : Application | Difficulté : Moyenne]} \\
    En utilisant \texttt{divmod\_poly}, écrire \texttt{pgcd\_poly(A, B)} qui retourne le PGCD de $A$ et $B$ (unitaire) par divisions successives.
    \item \textbf{[Bloom : Évaluation | Difficulté : Moyenne]} \\
    Tester $\mathrm{pgcd}(X^4 - 1, X^3 - X) = X^2 - 1$.
    \item \textbf{[Bloom : Synthèse | Difficulté : Difficile]} \\
    Vérifier numériquement la propriété :
    \[ \mathrm{pgcd}(X^m - 1, X^n - 1) = X^{\mathrm{pgcd}(m, n)} - 1 \]
    pour $10$ couples $(m, n)$ aléatoires avec $m, n \le 20$.
\end{enumerate}

\vspace{0.3cm}

\noindent\textbf{Exercice 5 : Bézout étendu pour polynômes}
\begin{enumerate}
    \item \textbf{[Bloom : Synthèse | Difficulté : Difficile]} \\
    Étendre l'algorithme d'Euclide en \texttt{bezout\_poly(A, B)} qui retourne $(D, U, V)$ tel que $D = A \wedge B$ et $A U + B V = D$. (Adapter l'algorithme d'Euclide étendu de l'arithmétique entière.)
    \item \textbf{[Bloom : Évaluation | Difficulté : Moyenne]} \\
    Tester sur $A = X^3 + 1$ et $B = X^2 + X$ : vérifier que $A \cdot 1 + B \cdot (1 - X) = X + 1$ (le PGCD).
    \item \textbf{[Bloom : Évaluation | Difficulté : Moyenne]} \\
    Comparer avec \texttt{sympy.gcdex(A, B, X)}.
\end{enumerate}

\vspace{0.3cm}

\noindent\textbf{Exercice 6 : Recherche de racines rationnelles}
\begin{enumerate}
    \item \textbf{[Bloom : Compréhension | Difficulté : Moyenne]} \\
    \emph{Théorème des racines rationnelles :} si $P = a_n X^n + \cdots + a_0 \in \mathbb{Z}[X]$ admet une racine rationnelle $p/q$ écrite sous forme irréductible, alors $p \mid a_0$ et $q \mid a_n$.
    Écrire \texttt{racines\_rationnelles(P)} qui retourne la liste de toutes les racines rationnelles d'un polynôme à coefficients entiers, en énumérant les couples $(p, q)$ candidats.
    \item \textbf{[Bloom : Évaluation | Difficulté : Moyenne]} \\
    Tester sur $P = 6X^3 - 11X^2 + 6X - 1$ (racines $1, 1/2, 1/3$).
    \item \textbf{[Bloom : Création | Difficulté : Difficile]} \\
    Améliorer en factorisant $P$ par chaque racine trouvée (Hörner) et en répétant la recherche sur le quotient.
\end{enumerate}

\vspace{0.3cm}

\noindent\textbf{Exercice 7 : Multiplicité d'une racine par dérivation}
\begin{enumerate}
    \item \textbf{[Bloom : Application | Difficulté : Moyenne]} \\
    Ajouter à votre classe \texttt{Polynome} la méthode \texttt{derivee(self)} qui retourne $P'$ (formule formelle).
    \item \textbf{[Bloom : Synthèse | Difficulté : Moyenne]} \\
    Écrire \texttt{multiplicite(P, a)} qui retourne la multiplicité de $a$ comme racine de $P$, en utilisant la caractérisation $P(a) = P'(a) = \cdots = P^{(m-1)}(a) = 0$ et $P^{(m)}(a) \neq 0$.
    \item \textbf{[Bloom : Évaluation | Difficulté : Moyenne]} \\
    Tester sur $P = (X - 1)^3 (X + 2)^2 (X - 5)$ développé : vérifier $\mathrm{mult}_1 = 3$, $\mathrm{mult}_{-2} = 2$, $\mathrm{mult}_5 = 1$.
\end{enumerate}

\vspace{0.3cm}

\noindent\textbf{Exercice 8 : Application --- Codes correcteurs Reed-Solomon (encodage)}
\begin{enumerate}
    \item \textbf{[Bloom : Compréhension | Difficulté : Difficile]} \\
    Le principe simple de Reed-Solomon : on encode un message $(m_0, m_1, \dots, m_{k-1}) \in \mathbb{F}_q^k$ comme le polynôme $M(X) = \sum m_i X^i$, puis on l'évalue en $n$ points distincts $\alpha_1, \dots, \alpha_n \in \mathbb{F}_q$ (avec $n > k$). Le mot de code est $(M(\alpha_1), \dots, M(\alpha_n))$.

    Écrire \texttt{encoder\_RS(message, points)} en travaillant dans $\mathbb{Z}/p\mathbb{Z}$ avec $p$ premier (par exemple $p = 11$).
    \item \textbf{[Bloom : Évaluation | Difficulté : Moyenne]} \\
    Encoder le message $(3, 1, 4)$ avec les points d'évaluation $\{1, 2, 3, 4, 5\}$ dans $\mathbb{F}_{11}$. Donner le mot de code obtenu.
    \item \textbf{[Bloom : Synthèse | Difficulté : Difficile]} \\
    \emph{Décodage simple (sans erreur)} : étant donné un mot de code $(y_1, \dots, y_n)$ supposé sans erreur et les points $(\alpha_1, \dots, \alpha_n)$, on retrouve le polynôme $M$ par interpolation de Lagrange. Implémenter \texttt{decoder\_RS\_sans\_erreur(mot, points, k)} qui retourne le message original.
\end{enumerate}

\chapter{Fractions rationnelles}

Ce chapitre construit le corps $K(X)$ des fractions rationnelles, puis en étudie l'arithmétique fine : forme irréductible, pôles, zéros, multiplicités, et enfin la \emph{décomposition en éléments simples}, outil central pour le calcul intégral, le calcul de séries télescopiques, le filtrage en traitement du signal et l'analyse des systèmes linéaires en automatique. Ce chapitre est l'aboutissement naturel du chapitre V : on quitte l'anneau $K[X]$ pour son \emph{corps des fractions}, par analogie au passage de $\mathbb{Z}$ à $\mathbb{Q}$.

\section{Corps des fractions $K(X)$}

\subsection{Motivation}
L'anneau $K[X]$ est intègre mais n'est \emph{pas} un corps : seuls les polynômes constants non nuls sont inversibles. On ne peut pas diviser librement par n'importe quel polynôme non nul. C'est exactement la situation rencontrée dans $\mathbb{Z}$ : pour pouvoir diviser, on a construit $\mathbb{Q}$, le corps des fractions de $\mathbb{Z}$. La même construction appliquée à $K[X]$ donne le corps $K(X)$ des \textbf{fractions rationnelles}, où l'expression $1/(X^2 + 1)$ devient un objet algébrique parfaitement défini.

Le passage à $K(X)$ ouvre la porte au calcul des primitives de fractions rationnelles, à l'analyse des transformées de Laplace en automatique, et plus généralement à toute l'analyse réelle et complexe d'expressions algébriques.

\subsection{Approche par problème : à quelles conditions $A/B = C/D$ ?}
\textbf{Problème :} Quand peut-on dire que $\frac{X - 1}{X^2 - 1}$ et $\frac{1}{X + 1}$ représentent la même fraction rationnelle ? Plus généralement, quelle relation $\sim$ identifier sur les couples $(A, B)$ avec $B \neq 0$ pour obtenir un cadre algébrique cohérent ?

\textbf{Résolution :} Dans $\mathbb{Q}$, $\frac{a}{b} = \frac{c}{d} \iff ad = bc$. La même règle s'applique aux polynômes :
\[ \frac{X - 1}{X^2 - 1} = \frac{1}{X + 1} \iff (X - 1)(X + 1) = (X^2 - 1) \cdot 1, \]
ce qui est vrai. Donc on définit $(A, B) \sim (C, D) \iff AD = BC$ et l'on note $A/B$ la classe d'équivalence de $(A, B)$. Cette classe d'équivalence est notre \emph{fraction rationnelle}.

\subsection{Construction de $K(X)$}

\begin{definition}[Fraction rationnelle, ensemble $K(X)$]
Sur l'ensemble $K[X] \times (K[X] \setminus \{0\})$, on définit la relation $\sim$ par :
\[ (A, B) \sim (C, D) \iff AD = BC. \]
C'est une relation d'équivalence (vérification directe, exploitant l'intégrité de $K[X]$).
L'ensemble \textbf{des fractions rationnelles} à coefficients dans $K$ est l'ensemble quotient
\[ K(X) = \big( K[X] \times (K[X] \setminus \{0\}) \big) \big/ \sim, \]
et l'on note la classe de $(A, B)$ par $\dfrac{A}{B}$.
\end{definition}

\begin{remarque}
La condition $B \neq 0$ est essentielle : on ne peut pas diviser par le polynôme nul, exactement comme dans $\mathbb{Q}$.
\end{remarque}

\subsection{Opérations sur $K(X)$}

\begin{definition}[Addition, multiplication, opposé, inverse]
Soient $\dfrac{A}{B}$ et $\dfrac{C}{D}$ deux fractions de $K(X)$. On pose :
\begin{align*}
\frac{A}{B} + \frac{C}{D} &= \frac{AD + BC}{BD}, \\
\frac{A}{B} \cdot \frac{C}{D} &= \frac{AC}{BD}, \\
-\frac{A}{B} &= \frac{-A}{B}.
\end{align*}
Si de plus $A \neq 0$, on définit l'inverse multiplicatif :
\[ \left(\frac{A}{B}\right)^{-1} = \frac{B}{A}. \]
\end{definition}

\begin{proposition}[Bonne définition]
Les opérations ci-dessus sont indépendantes du choix des représentants des classes : si $\frac{A}{B} = \frac{A'}{B'}$ et $\frac{C}{D} = \frac{C'}{D'}$, alors $\frac{AD + BC}{BD} = \frac{A'D' + B'C'}{B'D'}$ et de même pour le produit.
\end{proposition}

\begin{proof}
Vérification par calcul direct utilisant les égalités $AB' = A'B$ et $CD' = C'D$ et l'intégrité de $K[X]$.
\end{proof}

\subsection{Structure de corps}

\begin{theoreme}[$K(X)$ est un corps commutatif]
$(K(X), +, \times)$ est un \emph{corps commutatif} :
\begin{itemize}
    \item l'élément neutre additif est $\dfrac{0}{1}$, noté $0$ ;
    \item l'élément neutre multiplicatif est $\dfrac{1}{1}$, noté $1$ ;
    \item toute fraction non nulle $\dfrac{A}{B}$ admet pour inverse $\dfrac{B}{A}$.
\end{itemize}
\end{theoreme}

\begin{proof}
Les axiomes d'anneau commutatif (associativité, commutativité, distributivité) se vérifient directement par calcul sur les représentants. L'inverse de $\dfrac{A}{B}$ est $\dfrac{B}{A}$ lorsque $A \neq 0$, car
\[ \frac{A}{B} \cdot \frac{B}{A} = \frac{AB}{BA} = \frac{AB}{AB} = \frac{1}{1} = 1. \]
\end{proof}

\subsection{Inclusion $K[X] \hookrightarrow K(X)$}

\begin{proposition}[{Plongement de $K[X]$ dans $K(X)$}]
L'application $j : K[X] \to K(X),\ P \mapsto \dfrac{P}{1}$ est un morphisme d'anneaux \emph{injectif}. On identifie $K[X]$ à son image dans $K(X)$ : tout polynôme est une fraction rationnelle particulière (de dénominateur $1$).
\end{proposition}

\begin{proof}
$j$ est un morphisme : $j(P + Q) = \frac{P + Q}{1} = \frac{P}{1} + \frac{Q}{1} = j(P) + j(Q)$ ; idem pour le produit ; $j(1) = 1$.
Injectivité : $j(P) = 0 \iff \frac{P}{1} = \frac{0}{1} \iff P = 0$.
\end{proof}

\begin{remarque}[Universalité du corps des fractions]
$K(X)$ est le \emph{plus petit} corps contenant $K[X]$ : si $L$ est un corps et $\varphi : K[X] \to L$ un morphisme injectif d'anneaux, alors $\varphi$ se prolonge de manière unique en un morphisme injectif $K(X) \to L$. Cette propriété universelle caractérise $K(X)$ à isomorphisme près.
\end{remarque}

\subsection{Cas particuliers : $\mathbb{R}(X)$ et $\mathbb{C}(X)$}

\begin{itemize}
    \item $\mathbb{R}(X)$ est le corps des fractions rationnelles à coefficients réels.
    \item $\mathbb{C}(X)$ est le corps des fractions rationnelles à coefficients complexes.
    \item L'inclusion $\mathbb{R}[X] \subset \mathbb{C}[X]$ induit $\mathbb{R}(X) \subset \mathbb{C}(X)$.
    \item La conjugaison complexe se prolonge à $\mathbb{C}(X)$ : si $F = A/B$, on pose $\bar{F} = \bar{A}/\bar{B}$ ; c'est un automorphisme de corps de $\mathbb{C}(X)$, dont les invariants sont exactement $\mathbb{R}(X)$.
\end{itemize}

\subsection{Fonction rationnelle associée}

\begin{definition}[Fonction rationnelle]
Soit $F = A/B \in K(X)$. La \textbf{fonction rationnelle} associée est l'application
\[ \widetilde{F} : K \setminus Z(B) \longrightarrow K, \qquad x \longmapsto \frac{A(x)}{B(x)}, \]
où $Z(B) = \{x \in K : B(x) = 0\}$ est l'ensemble des racines (zéros) de $B$.
\end{definition}

\begin{remarque}
Un même $F \in K(X)$ peut s'écrire avec différents représentants $A/B$. Le domaine de définition de $\widetilde{F}$ peut donc varier selon l'écriture choisie. C'est pourquoi on travaille avec la \emph{forme irréductible} (§2), où le domaine est uniquement déterminé par la fraction rationnelle.
\end{remarque}

\subsection{Exemples résolus (Difficulté ascendante)}

\textbf{Exemple 1 (Simplification --- Facile).}
Vérifier que $\dfrac{X^2 - 1}{X^2 - 2X + 1} = \dfrac{X + 1}{X - 1}$ dans $\mathbb{R}(X)$.

\textit{Correction :} on factorise numérateur et dénominateur :
$X^2 - 1 = (X - 1)(X + 1)$ et $X^2 - 2X + 1 = (X - 1)^2$.
$\dfrac{(X-1)(X+1)}{(X-1)^2} = \dfrac{X + 1}{X - 1}$. \checkmark

(Vérification par produit en croix : $(X^2 - 1)(X - 1) = (X - 1)(X + 1)(X - 1) = (X^2 - 2X + 1)(X + 1)$. \checkmark)

\medskip

\textbf{Exemple 2 (Addition de fractions --- Intermédiaire).}
Calculer $\dfrac{1}{X} + \dfrac{1}{X + 1} - \dfrac{1}{X^2 + X}$.

\textit{Correction :} le dénominateur commun est $X(X + 1) = X^2 + X$.
\[ \frac{1}{X} + \frac{1}{X+1} - \frac{1}{X^2 + X} = \frac{(X + 1) + X - 1}{X(X + 1)} = \frac{2X}{X(X+1)} = \frac{2}{X + 1}. \]

\medskip

\textbf{Exemple 3 (Inverse et calcul --- Difficile).}
Soit $F = \dfrac{X^2 + 1}{X - 2}$. Calculer $F + F^{-1}$ et donner le résultat sous forme $A/B$ avec $A, B \in \mathbb{R}[X]$.

\textit{Correction :} $F^{-1} = \dfrac{X - 2}{X^2 + 1}$.
\[ F + F^{-1} = \frac{X^2 + 1}{X - 2} + \frac{X - 2}{X^2 + 1} = \frac{(X^2 + 1)^2 + (X - 2)^2}{(X - 2)(X^2 + 1)}. \]
Développement du numérateur : $(X^2 + 1)^2 + (X - 2)^2 = X^4 + 2X^2 + 1 + X^2 - 4X + 4 = X^4 + 3X^2 - 4X + 5$.
Développement du dénominateur : $(X - 2)(X^2 + 1) = X^3 - 2X^2 + X - 2$.
\[ F + F^{-1} = \frac{X^4 + 3X^2 - 4X + 5}{X^3 - 2X^2 + X - 2}. \]

\subsection{Exercices pratiques avec Python}
\texttt{sympy} traite nativement les fractions rationnelles : addition, multiplication, simplification, mise au même dénominateur via \texttt{together} et \texttt{cancel}.

\begin{lstlisting}
import sympy as sp

X = sp.symbols('X')

# 1. Operations elementaires
F1 = (X**2 + 1) / (X - 2)
F2 = (X - 2) / (X**2 + 1)
print(f"F1        = {F1}")
print(f"F2 = 1/F1 = {F2}")
print(f"F1 * F2   = {sp.simplify(F1 * F2)}    (attendu : 1)")

# 2. Addition et mise au meme denominateur
G = 1/X + 1/(X + 1) - 1/(X**2 + X)
print(f"\n1/X + 1/(X+1) - 1/(X^2+X) = {sp.together(G)}")
print(f"  simplifie : {sp.simplify(G)}")

# 3. F + 1/F : verification de l'exemple 3
H = F1 + F2
print(f"\nF + 1/F = {sp.together(H)}")
print(f"  simplifie : {sp.simplify(H)}")

# 4. Test : deux fractions sont-elles egales ?
A1 = (X**2 - 1) / (X**2 - 2*X + 1)
A2 = (X + 1) / (X - 1)
print(f"\n(X^2-1)/(X^2-2X+1) - (X+1)/(X-1) = {sp.simplify(A1 - A2)}")
\end{lstlisting}

\subsection{Exercices à faire}
\begin{enumerate}
    \item \textbf{Exercice 1.} Vérifier que $\dfrac{X^3 - 1}{X - 1} = X^2 + X + 1$ dans $\mathbb{R}(X)$, en explicitant le sens donné à cette égalité.
    \item \textbf{Exercice 2.} Effectuer la somme $\dfrac{1}{X - 1} + \dfrac{1}{X + 1} + \dfrac{2}{X^2 - 1}$ et la simplifier au maximum.
    \item \textbf{Exercice 3.} Soit $F = \dfrac{X^2 + 2X + 1}{X^3 + X^2}$ dans $\mathbb{R}(X)$. Donner deux écritures de $F$ avec des dénominateurs différents.
    \item \textbf{Exercice 4.} Démontrer que la relation $\sim$ définie sur $K[X] \times (K[X] \setminus \{0\})$ par $(A, B) \sim (C, D) \iff AD = BC$ est bien une relation d'équivalence. Où utilise-t-on l'intégrité de $K[X]$ ?
    \item \textbf{Exercice 5 (synthèse).} Soit $F \in \mathbb{C}(X)$. Démontrer que $F \in \mathbb{R}(X)$ si et seulement si $\bar{F} = F$, où $\bar{F}$ est obtenu en conjuguant les coefficients d'un représentant de $F$. (Indication : commencer par vérifier que la conjugaison est bien définie sur les classes.)
\end{enumerate}

\section{Forme irréductible et fonctions rationnelles}

\subsection{Motivation}
Une même fraction rationnelle $F \in K(X)$ admet une infinité d'écritures $A/B$. Pour calculer, factoriser, étudier des limites ou simplifier des expressions, on a besoin d'une \emph{écriture canonique}, libre de tout facteur commun parasite : c'est la \textbf{forme irréductible}. Elle est l'analogue de la forme « fraction réduite » d'un rationnel ($\frac{6}{8} = \frac{3}{4}$). Au-delà de la commodité, la forme irréductible permet de définir intrinsèquement le domaine de définition d'une fonction rationnelle, ses pôles, ses zéros (étudiés en §3).

\subsection{Approche par problème : trouver la « plus simple » écriture}
\textbf{Problème :} Réduire la fraction rationnelle $F = \dfrac{X^4 - 1}{X^3 - X^2 + X - 1}$ à sa forme la plus simple.

\textbf{Résolution :} factorisons numérateur et dénominateur.
$X^4 - 1 = (X^2 - 1)(X^2 + 1) = (X - 1)(X + 1)(X^2 + 1)$.
$X^3 - X^2 + X - 1 = X^2(X - 1) + (X - 1) = (X - 1)(X^2 + 1)$.
Le PGCD est $(X - 1)(X^2 + 1)$. En simplifiant :
\[ F = \frac{(X - 1)(X + 1)(X^2 + 1)}{(X - 1)(X^2 + 1)} = X + 1. \]
Cette écriture est minimale : aucun facteur commun ne subsiste entre numérateur ($X + 1$) et dénominateur ($1$). Elle s'obtient algorithmiquement en divisant numérateur et dénominateur par leur PGCD.

\subsection{Forme irréductible}

\begin{definition}[Forme irréductible]
Une fraction rationnelle $F = \dfrac{A}{B} \in K(X)$ est dite sous \textbf{forme irréductible} si $A \wedge B = 1$, c'est-à-dire si $A$ et $B$ sont premiers entre eux dans $K[X]$.

Une forme irréductible $A/B$ est dite \emph{normalisée} si de plus $B$ est unitaire (coefficient dominant égal à $1$).
\end{definition}

\begin{theoreme}[Existence et unicité de la forme irréductible]
Toute fraction rationnelle $F \in K(X)$, $F \neq 0$, admet :
\begin{enumerate}
    \item au moins une forme irréductible ;
    \item une unique forme irréductible \emph{normalisée}.
\end{enumerate}
\end{theoreme}

\begin{proof}
\emph{Existence.} Soit $A/B$ un représentant quelconque de $F$. Notons $D = A \wedge B$. Alors $A = D A'$ et $B = D B'$ avec $A' \wedge B' = 1$. Donc $\dfrac{A}{B} = \dfrac{A'}{B'}$ est une forme irréductible.

\emph{Unicité de la forme normalisée.} Supposons $\dfrac{A_1}{B_1} = \dfrac{A_2}{B_2}$ avec $A_i \wedge B_i = 1$ et $B_i$ unitaire. L'égalité $A_1 B_2 = A_2 B_1$ et le théorème de Gauss donnent : comme $B_1 \mid A_1 B_2$ et $B_1 \wedge A_1 = 1$, alors $B_1 \mid B_2$. De même $B_2 \mid B_1$. Comme $B_1, B_2$ sont unitaires, $B_1 = B_2$. Ensuite $A_1 = A_2$.
\end{proof}

\begin{remarque}
Sans la condition de normalisation, la forme irréductible n'est unique qu'à \emph{multiplication par une constante non nulle près} : $A/B$ et $(\lambda A)/(\lambda B)$ sont deux formes irréductibles équivalentes pour tout $\lambda \in K^*$.
\end{remarque}

\subsection{Algorithme de réduction}

L'algorithme suivant, basé sur l'algorithme d'Euclide vu au chapitre V, calcule la forme irréductible normalisée d'une fraction $A/B$ donnée :

\begin{enumerate}
    \item Calculer $D = A \wedge B$ par l'algorithme d'Euclide.
    \item Effectuer les divisions : $A' = A / D$, $B' = B / D$.
    \item Diviser $A'$ et $B'$ par le coefficient dominant de $B'$ pour rendre $B'$ unitaire.
\end{enumerate}

Le résultat $A''/B''$ est l'unique forme irréductible normalisée de $F$.

\subsection{Domaine de la fonction rationnelle}

\begin{proposition}[Domaine intrinsèque]
Soit $F \in K(X)$ et $A/B$ sa forme irréductible. La fonction rationnelle $\widetilde{F} : x \mapsto A(x)/B(x)$ est définie sur $K \setminus Z(B)$, où $Z(B)$ est l'ensemble des racines de $B$. Cet ensemble \emph{ne dépend pas} du choix de la forme irréductible : il est intrinsèque à $F$.
\end{proposition}

\begin{proof}
Si $A/B$ et $A'/B'$ sont deux formes irréductibles de la même fraction $F$, alors $B$ et $B'$ sont associés (par l'unicité ci-dessus à un scalaire près). Donc $Z(B) = Z(B')$.
\end{proof}

\begin{definition}[Pôle, zéro --- définitions provisoires]
Soit $F \in K(X)$ donnée par sa forme irréductible $A/B$.
\begin{itemize}
    \item Un \textbf{zéro} de $F$ est une racine de $A$ dans $K$ (ou son extension $\mathbb{C}$).
    \item Un \textbf{pôle} de $F$ est une racine de $B$ dans $K$ (ou son extension $\mathbb{C}$).
\end{itemize}
Les multiplicités de zéros et pôles seront formalisées au §3.
\end{definition}

\begin{remarque}
La nécessité de la forme irréductible apparaît clairement ici : sans simplification, la fraction $\dfrac{X(X-1)}{X-1}$ semblerait avoir un pôle en $1$, mais une fois réduite à $X$ on voit qu'il n'en est rien. Les pôles et zéros sont des invariants algébriques de $F$, lisibles seulement sur la forme irréductible.
\end{remarque}

\subsection{Opérations sur les formes irréductibles}

Les opérations $+, \times, /$ ne préservent pas la forme irréductible : la somme de deux fractions sous forme irréductible doit en général être réduite à nouveau. Toutefois :

\begin{proposition}[Multiplication et inverse]
Soient $F = A/B$ et $G = C/D$ deux fractions rationnelles non nulles sous forme irréductible.
\begin{itemize}
    \item Si $A \wedge D = 1$ et $C \wedge B = 1$, alors $\dfrac{AC}{BD}$ est sous forme irréductible.
    \item L'inverse de $A/B$ est $B/A$, et il est sous forme irréductible (puisque $A \wedge B = 1$).
\end{itemize}
\end{proposition}

\begin{proof}
Pour la multiplication : un facteur commun à $AC$ et $BD$ devrait, par Gauss, diviser un des facteurs ; les hypothèses excluent toutes les combinaisons non triviales. Pour l'inverse : la condition $B \wedge A = 1$ est symétrique.
\end{proof}

\subsection{Exemples résolus (Difficulté ascendante)}

\textbf{Exemple 1 (Réduction directe --- Facile).}
Mettre sous forme irréductible $F = \dfrac{2X^2 - 8}{4X - 8}$ dans $\mathbb{R}(X)$.

\textit{Correction :} on factorise $2X^2 - 8 = 2(X - 2)(X + 2)$ et $4X - 8 = 4(X - 2)$.
\[ F = \frac{2(X - 2)(X + 2)}{4(X - 2)} = \frac{2(X + 2)}{4} = \frac{X + 2}{2}. \]
La forme normalisée demande un dénominateur unitaire : ici $2$ n'est pas unitaire, mais on accepte $1/2 \cdot (X + 2)$ ou bien on écrit $F = \dfrac{X/2 + 1}{1}$ : c'est en fait un polynôme.

(Remarque : les fractions « polynomiales » ont pour forme irréductible normalisée $P/1$ avec $P$ polynôme.)

\medskip

\textbf{Exemple 2 (Réduction par algorithme d'Euclide --- Intermédiaire).}
Mettre sous forme irréductible $F = \dfrac{X^3 - X}{X^2 - 1}$ dans $\mathbb{R}(X)$.

\textit{Correction :}
$A = X^3 - X = X(X^2 - 1)$ et $B = X^2 - 1$. Par Euclide : $A = X \cdot B + 0$, donc $A \wedge B = B = X^2 - 1$.
$A / D = X(X^2 - 1)/(X^2 - 1) = X$.
$B / D = (X^2 - 1)/(X^2 - 1) = 1$.
Forme irréductible normalisée : $F = X$.

\medskip

\textbf{Exemple 3 (Domaine de définition --- Difficile).}
Donner la forme irréductible et le domaine de définition de la fonction $x \mapsto \dfrac{x^4 - 5x^2 + 4}{x^4 - 1}$ sur $\mathbb{R}$.

\textit{Correction :} factorisons.
$X^4 - 5X^2 + 4 = (X^2 - 1)(X^2 - 4) = (X - 1)(X + 1)(X - 2)(X + 2)$.
$X^4 - 1 = (X^2 - 1)(X^2 + 1) = (X - 1)(X + 1)(X^2 + 1)$.
PGCD $= (X - 1)(X + 1) = X^2 - 1$.
Après simplification :
\[ F = \frac{(X - 2)(X + 2)}{X^2 + 1} = \frac{X^2 - 4}{X^2 + 1}. \]
Comme $X^2 + 1$ n'a pas de racine réelle, le domaine de définition de $\widetilde{F}$ sur $\mathbb{R}$ est $\mathbb{R}$ tout entier.

\smallskip

\textbf{Attention à un piège fréquent} : la fonction $x \mapsto \dfrac{x^4 - 5x^2 + 4}{x^4 - 1}$ telle qu'écrite n'est pas définie en $x = \pm 1$. La forme irréductible donne une fonction \emph{prolongée par continuité} en ces points (avec valeur $-3/2$ par exemple en $x = 1$). C'est précisément l'intérêt de la forme irréductible : elle révèle les vrais pôles, en éliminant les « fausses singularités » dues à des facteurs communs.

\subsection{Exercices pratiques avec Python}
\texttt{sympy.cancel} renvoie la forme irréductible d'une fraction rationnelle.

\begin{lstlisting}
import sympy as sp

X = sp.symbols('X')

# 1. Reduction explicite
F = (X**3 - X) / (X**2 - 1)
print(f"F            = {F}")
print(f"  forme irreductible : {sp.cancel(F)}")

# 2. Algorithme manuel : pgcd numerateur/denominateur
def forme_irreductible(num, den, var):
    g = sp.gcd(num, den)
    return sp.simplify(num / g), sp.simplify(den / g)

num = X**4 - 5*X**2 + 4
den = X**4 - 1
A, B = forme_irreductible(num, den, X)
print(f"\n{num} / {den}")
print(f"  pgcd     = {sp.gcd(num, den)}")
print(f"  forme irred. normalisee : ({A}) / ({B})")

# 3. Domaine reel d'une fonction rationnelle
def domaine_reel(num, den, var):
    """Retourne les valeurs reelles a exclure (poles)."""
    A, B = forme_irreductible(num, den, var)
    return sp.solve(B, var)

print(f"\nDomaine de F = ({num})/({den}) :")
print(f"  poles reels (apres simplification) : {domaine_reel(num, den, X)}")
print(f"  poles apparents si on n'avait pas simplifie : {sp.solve(den, X)}")
\end{lstlisting}

\subsection{Exercices à faire}
\begin{enumerate}
    \item \textbf{Exercice 1.} Mettre sous forme irréductible normalisée chacune des fractions suivantes dans $\mathbb{R}(X)$ :
    \begin{itemize}
        \item $F_1 = \dfrac{3X^2 - 12}{6X + 12}$ ;
        \item $F_2 = \dfrac{X^4 + X^3}{X^3 + X^2}$ ;
        \item $F_3 = \dfrac{X^3 - 1}{X^2 + X + 1}$.
    \end{itemize}
    \item \textbf{Exercice 2.} Soient $F = \dfrac{X^2 + 1}{X - 1}$ et $G = \dfrac{X^2 - 1}{X + 2}$ dans $\mathbb{R}(X)$, sous forme irréductible. Calculer $F + G$ et donner le résultat sous forme irréductible normalisée.
    \item \textbf{Exercice 3.} Soit $F = \dfrac{X^2 - 4}{X^3 - 2X^2 - X + 2}$. Mettre $F$ sous forme irréductible et déterminer son domaine de définition sur $\mathbb{R}$.
    \item \textbf{Exercice 4.} Démontrer que si $F = A/B$ est sous forme irréductible avec $\deg A < \deg B$, alors la même fraction multipliée par tout polynôme non nul $P$ (i.e.\ $F \cdot P$) reste \emph{strictement propre} (degré du numérateur strictement inférieur à celui du dénominateur) si et seulement si $\deg P < \deg B - \deg A$.
    \item \textbf{Exercice 5 (synthèse).} Soit $F \in \mathbb{R}(X)$ sous forme irréductible $A/B$. Montrer que $F$ et $-F$ ont la même forme irréductible normalisée à un signe près sur le numérateur. Que dire de $F$ et $1/F$ (lorsque $F \neq 0$) ?
\end{enumerate}

\section{Degré, pôles, zéros, multiplicités}

\subsection{Motivation}
Une fraction rationnelle est entièrement caractérisée localement par trois invariants : son \emph{degré} (qui contrôle son comportement à l'infini), la \emph{multiplicité de ses pôles} (qui contrôle l'explosion près des points singuliers) et la \emph{multiplicité de ses zéros} (qui contrôle la nullité). Ces invariants sont à la base de toute analyse de fractions rationnelles : calcul de primitives ($\int dx/(x-a)^m$), étude asymptotique en automatique ($1/(p^2 + 2\zeta\omega p + \omega^2)$), théorème des résidus en analyse complexe, et finalement la décomposition en éléments simples (§4).

\subsection{Approche par problème}
\textbf{Problème :} On considère $F = \dfrac{X^3 - X}{(X - 1)^2 (X^2 + 4)}$ dans $\mathbb{R}(X)$, sous forme \emph{déjà irréductible} (à vérifier).
\begin{enumerate}
    \item Quel est le comportement de $\widetilde{F}(x)$ lorsque $x \to +\infty$ ?
    \item Combien de pôles réels ? Avec quelle multiplicité chacun ?
    \item Combien de zéros réels ? Avec quelle multiplicité ?
\end{enumerate}

\textbf{Résolution :}
$\deg(X^3 - X) = 3$, $\deg((X-1)^2 (X^2 + 4)) = 4$. Donc $\deg F = 3 - 4 = -1$ : à l'infini, $F(x) \sim c/x$ et tend vers $0$.
Le numérateur factorisé : $X^3 - X = X(X - 1)(X + 1)$. PGCD avec $(X - 1)^2 (X^2 + 4)$ : facteur commun $X - 1$. Donc $F$ \emph{n'était pas} sous forme irréductible.
Forme irréductible : $F = \dfrac{X(X + 1)}{(X - 1)(X^2 + 4)}$.
Pôles réels : $X - 1 = 0$, soit $x = 1$, multiplicité $1$. Le facteur $X^2 + 4$ n'a pas de racine réelle (mais a deux pôles complexes $\pm 2i$ de multiplicité $1$).
Zéros réels : $X(X + 1) = 0$, soit $x = 0$ (multiplicité $1$) et $x = -1$ (multiplicité $1$).

Cette section formalise ces calculs.

\subsection{Degré d'une fraction rationnelle}

\begin{definition}[Degré]
Soit $F = A/B \in K(X)$ avec $A \neq 0$ et $B \neq 0$. Le \textbf{degré} de $F$ est :
\[ \deg F = \deg A - \deg B \in \mathbb{Z}. \]
Par convention, $\deg 0 = -\infty$.
\end{definition}

\begin{proposition}[Bonne définition]
$\deg F$ ne dépend pas du représentant choisi de $F$.
\end{proposition}

\begin{proof}
Si $A/B = A'/B'$, alors $A B' = A' B$, donc $\deg A + \deg B' = \deg A' + \deg B$, soit $\deg A - \deg B = \deg A' - \deg B'$.
\end{proof}

\begin{proposition}[Propriétés du degré]
Pour $F, G \in K(X)$ :
\begin{enumerate}
    \item $\deg(FG) = \deg F + \deg G$ ;
    \item $\deg(F + G) \le \max(\deg F, \deg G)$, avec égalité si $\deg F \neq \deg G$ ;
    \item $\deg(1/F) = -\deg F$ pour $F \neq 0$ ;
    \item Si $F \in K[X] \subset K(X)$, le degré comme fraction coïncide avec le degré comme polynôme.
\end{enumerate}
\end{proposition}

\subsection{Fractions propres et partie entière}

\begin{definition}[Fraction propre, strictement propre]
$F = A/B \in K(X)$ est dite :
\begin{itemize}
    \item \textbf{propre} si $\deg F \le 0$ (i.e.\ $\deg A \le \deg B$) ;
    \item \textbf{strictement propre} si $\deg F < 0$ (i.e.\ $\deg A < \deg B$) ;
    \item \textbf{impropre} si $\deg F > 0$.
\end{itemize}
\end{definition}

\begin{theoreme}[Décomposition partie entière + reste strictement propre]
Pour toute $F = A/B \in K(X)$ avec $B \neq 0$, il existe un \emph{unique} couple $(E, R)$ tel que :
\[ F = E + \frac{R}{B}, \qquad E \in K[X], \qquad \deg R < \deg B. \]
$E$ est appelée la \textbf{partie entière} (ou partie polynomiale) de $F$, et $R/B$ est la \emph{partie strictement propre}. Lorsque $F$ est strictement propre, $E = 0$.
\end{theoreme}

\begin{proof}
Existence : la division euclidienne dans $K[X]$ donne $A = B Q + R$ avec $\deg R < \deg B$. Alors $F = A/B = Q + R/B$, donc $E = Q$ convient.
Unicité : si $E + R/B = E' + R'/B$, alors $E - E' = (R' - R)/B$ ; le membre de gauche est polynomial, donc $B \mid R' - R$. Comme $\deg(R' - R) < \deg B$, on a $R' = R$, puis $E = E'$.
\end{proof}

\begin{remarque}
La partie entière $E$ donne le \emph{comportement asymptotique} : pour $|x| \to \infty$, $F(x) \sim E(x)$. C'est l'analogue du « polynôme dominant » en analyse réelle.
\end{remarque}

\subsection{Multiplicité d'un zéro et d'un pôle}

\begin{definition}[Multiplicité d'un zéro et d'un pôle]
Soit $F \in K(X)$ non nulle, sous forme irréductible $A/B$, et $a$ un élément de $K$ (ou plus généralement de $\mathbb{C}$ pour $F \in \mathbb{R}(X)$).
\begin{itemize}
    \item $a$ est un \textbf{zéro} de $F$ de multiplicité $m \ge 1$ si $(X - a)^m \mid A$ et $(X - a)^{m+1} \nmid A$. On note $\mathrm{mult}^0_a(F) = m$.
    \item $a$ est un \textbf{pôle} de $F$ de multiplicité $m \ge 1$ si $(X - a)^m \mid B$ et $(X - a)^{m+1} \nmid B$. On note $\mathrm{mult}^\infty_a(F) = m$.
\end{itemize}
Si $a$ n'est ni zéro ni pôle, on pose $\mathrm{mult}^0_a(F) = 0$ et $\mathrm{mult}^\infty_a(F) = 0$.
\end{definition}

\begin{remarque}
Comme $A \wedge B = 1$, un même $a$ ne peut pas être à la fois zéro et pôle : au plus l'un des deux est non nul. C'est précisément l'intérêt de la forme irréductible.
\end{remarque}

\subsection{Comportement local près d'un pôle}

\begin{proposition}[Asymptotique au voisinage d'un pôle]
Si $a \in K$ est un pôle de $F$ de multiplicité $m$, alors il existe une fraction $G \in K(X)$, définie en $a$ et telle que $G(a) \neq 0$, vérifiant
\[ F = \frac{G}{(X - a)^m}. \]
En particulier, lorsque $K = \mathbb{R}$ et $a$ est un pôle réel, on a $\widetilde{F}(x) \underset{x \to a}{\sim} \dfrac{G(a)}{(x - a)^m}$.
\end{proposition}

\begin{proof}
Sous forme irréductible $A/B$, on factorise $B = (X - a)^m B'$ avec $B'(a) \neq 0$, puis $G = A/B' \in K(X)$.
\end{proof}

\subsection{Cas $\mathbb{C}(X)$ : relation entre degré, zéros et pôles}

\begin{theoreme}[Bilan zéros--pôles dans $\mathbb{C}(X)$]
Soit $F \in \mathbb{C}(X)$ non nulle, sous forme irréductible $A/B$. La somme des multiplicités des zéros (dans $\mathbb{C}$) moins la somme des multiplicités des pôles (dans $\mathbb{C}$) est égale à $\deg F$ :
\[ \sum_{a \in \mathbb{C}} \mathrm{mult}^0_a(F) - \sum_{a \in \mathbb{C}} \mathrm{mult}^\infty_a(F) = \deg A - \deg B = \deg F. \]
\end{theoreme}

\begin{proof}
Sur $\mathbb{C}$, $A$ et $B$ sont scindés. La somme des multiplicités des zéros vaut $\deg A$, celle des pôles $\deg B$.
\end{proof}

\begin{corollaire}
Une fraction rationnelle non constante de $\mathbb{C}(X)$ admet (en comptant multiplicités) autant de zéros que de pôles si et seulement si $\deg F = 0$ (i.e.\ $\deg A = \deg B$).
\end{corollaire}

\subsection{Exemples résolus (Difficulté ascendante)}

\textbf{Exemple 1 (Degré et partie entière --- Facile).}
Soit $F = \dfrac{X^4 + 1}{X^2 - 1}$. Calculer $\deg F$ et donner la décomposition partie entière + reste.

\textit{Correction :} $\deg F = 4 - 2 = 2$, donc $F$ est impropre. Division euclidienne :
$X^4 + 1 = (X^2 - 1)(X^2 + 1) + 2$, donc
\[ F = X^2 + 1 + \frac{2}{X^2 - 1}. \]
Partie entière : $E = X^2 + 1$. Reste strictement propre : $\dfrac{2}{X^2 - 1}$.

\medskip

\textbf{Exemple 2 (Pôles et zéros avec multiplicités --- Intermédiaire).}
Soit $F = \dfrac{(X - 2)^3 (X + 1)}{(X - 1)^2 (X + 1)^2}$ dans $\mathbb{R}(X)$. Mettre sous forme irréductible et déterminer les pôles et zéros avec leurs multiplicités.

\textit{Correction :} on simplifie le facteur commun $(X + 1)$ :
\[ F = \frac{(X - 2)^3}{(X - 1)^2 (X + 1)}. \]
Cette écriture est irréductible (les facteurs $(X - 2)$, $(X - 1)$, $(X + 1)$ sont premiers entre eux deux à deux).
\begin{itemize}
    \item Zéro : $a = 2$ de multiplicité $3$.
    \item Pôles : $a = 1$ de multiplicité $2$ ; $a = -1$ de multiplicité $1$.
\end{itemize}
Vérification du bilan : $3 - (2 + 1) = 0 = \deg F = 3 - 3$. \checkmark

\medskip

\textbf{Exemple 3 (Comportement asymptotique --- Difficile).}
Étudier le comportement de $\widetilde{F}(x)$ près de chaque pôle réel et à l'infini, pour $F = \dfrac{X^2 + 1}{(X - 1)^3 (X + 2)}$ dans $\mathbb{R}(X)$ (forme déjà irréductible).

\textit{Correction :}
\begin{itemize}
    \item À l'infini : $\deg F = 2 - 4 = -2$, donc $\widetilde{F}(x) \underset{x \to \pm \infty}{\sim} \dfrac{1}{x^2}$. La fonction tend vers $0$ « comme $1/x^2$ ».
    \item Pôle $a = 1$ de multiplicité $3$ : on écrit $F = \dfrac{(X^2 + 1)/(X + 2)}{(X - 1)^3}$. Au voisinage de $1$, le numérateur vaut $G(1) = 2/3$. Donc $\widetilde{F}(x) \underset{x \to 1}{\sim} \dfrac{2/3}{(x - 1)^3}$. La fonction « explose » comme $\pm \infty$ selon le signe de $x - 1$.
    \item Pôle $a = -2$ de multiplicité $1$ : $F = \dfrac{(X^2 + 1)/(X - 1)^3}{X + 2}$. Numérateur en $-2$ vaut $5 / (-3)^3 = -5/27$. Donc $\widetilde{F}(x) \underset{x \to -2}{\sim} \dfrac{-5/27}{x + 2}$.
\end{itemize}

\subsection{Exercices pratiques avec Python}
\texttt{sympy} fournit des fonctions pour décomposer en partie entière + reste (\texttt{div} sur les polynômes), trouver les pôles (zéros du dénominateur après simplification) et leur multiplicité.

\begin{lstlisting}
import sympy as sp

X = sp.symbols('X')

# 1. Degre d'une fraction et partie entiere
F = (X**4 + 1) / (X**2 - 1)
F_simpl = sp.cancel(F)
num, den = sp.fraction(F_simpl)
print(f"F = {F}")
print(f"  deg F        = {sp.degree(num) - sp.degree(den)}")

# Partie entiere + reste : division euclidienne
Q, R = sp.div(num, den, X)
print(f"  partie entiere E = {Q}")
print(f"  reste R          = {R}")
print(f"  -> F = ({Q}) + ({R})/({den})")

# 2. Poles et zeros avec multiplicites
F = (X - 2)**3 * (X + 1) / ((X - 1)**2 * (X + 1)**2)
F_irred = sp.cancel(F)
num, den = sp.fraction(F_irred)
print(f"\nF (forme irred.) = {F_irred}")
print(f"  zeros (avec multiplicites)  : {sp.roots(num, X)}")
print(f"  poles (avec multiplicites)  : {sp.roots(den, X)}")

# 3. Bilan zeros/poles = deg F (sur C)
def bilan_zeros_poles(F, var):
    F_irr = sp.cancel(F)
    n, d = sp.fraction(F_irr)
    z = sum(sp.roots(n, var).values())
    p = sum(sp.roots(d, var).values())
    return z, p, z - p

F = (X**3 - 1) / (X**4 + X**2)
z, p, diff = bilan_zeros_poles(F, X)
print(f"\nF = {F}")
print(f"  somme multiplicites zeros : {z}")
print(f"  somme multiplicites poles : {p}")
print(f"  z - p = {diff}    (attendu : deg F = {sp.degree(sp.numer(sp.cancel(F))) - sp.degree(sp.denom(sp.cancel(F)))})")
\end{lstlisting}

\subsection{Exercices à faire}
\begin{enumerate}
    \item \textbf{Exercice 1.} Calculer $\deg F$ pour chacune des fractions :
    \begin{itemize}
        \item $F_1 = \dfrac{X^5 + 3X}{X^2 + 1}$ ;
        \item $F_2 = \dfrac{X + 1}{X^3 - X}$ ;
        \item $F_3 = \dfrac{X^4 - 1}{X^4 + 1}$.
    \end{itemize}
    Indiquer dans chaque cas si la fraction est propre, strictement propre, ou impropre.
    \item \textbf{Exercice 2.} Donner la partie entière et le reste strictement propre de $F = \dfrac{X^4 - 2X^3 + X - 1}{X^2 + X + 1}$.
    \item \textbf{Exercice 3.} Déterminer les pôles et zéros (avec multiplicités) dans $\mathbb{C}$ de chacune des fractions suivantes, après mise sous forme irréductible :
    \begin{itemize}
        \item $F_1 = \dfrac{X^3 - X}{X^2 - 1}$ ;
        \item $F_2 = \dfrac{X^2 + 1}{(X^2 - 1)^2}$ ;
        \item $F_3 = \dfrac{X^4 + 1}{X^4 - 1}$.
    \end{itemize}
    \item \textbf{Exercice 4.} Soit $F \in \mathbb{R}(X)$ non nulle. Démontrer que les pôles complexes non réels de $F$ apparaissent par couples conjugués, avec la même multiplicité.
    \item \textbf{Exercice 5 (synthèse).} Soit $F \in \mathbb{C}(X)$ avec $\deg F = -2$, ayant un unique pôle $a$ de multiplicité $2$. Montrer qu'il existe $\lambda \in \mathbb{C}^*$ tel que $F = \dfrac{\lambda}{(X - a)^2}$. Plus généralement, énoncer et démontrer le résultat pour $\deg F = -m$ avec un unique pôle de multiplicité $m$.
\end{enumerate}

\section{Décomposition en éléments simples}

\subsection{Motivation}
La \textbf{décomposition en éléments simples} (DES) écrit toute fraction rationnelle comme somme d'un polynôme et de termes élémentaires de la forme $\dfrac{1}{(X - a)^k}$ (et, sur $\mathbb{R}$, de termes $\dfrac{\alpha X + \beta}{(X^2 + pX + q)^k}$). C'est l'outil universel pour :
\begin{itemize}
    \item le \emph{calcul de primitives} de fractions rationnelles ($\int dx / (x - a)^k$, $\int dx / (x^2 + 1)$ se calculent immédiatement) ;
    \item l'\emph{inversion des transformées de Laplace et en $z$} en automatique et en traitement du signal ;
    \item le \emph{calcul de séries télescopiques} et de sommes infinies ;
    \item le \emph{théorème des résidus} en analyse complexe.
\end{itemize}
On démontre ici l'existence et l'unicité de la DES, puis l'on présente les méthodes de calcul des coefficients.

\subsection{Approche par problème : intégrale d'une fraction simple}
\textbf{Problème :} Calculer $\displaystyle \int_2^3 \frac{dx}{x^2 - 1}$.

\textbf{Résolution :} on écrit $\dfrac{1}{x^2 - 1} = \dfrac{1}{(x - 1)(x + 1)}$ et l'on cherche $\alpha, \beta \in \mathbb{R}$ tels que
\[ \frac{1}{(x - 1)(x + 1)} = \frac{\alpha}{x - 1} + \frac{\beta}{x + 1}. \]
Méthode de multiplication par $(x - 1)$ puis évaluation en $x = 1$ : $\alpha = \dfrac{1}{1 + 1} = \dfrac{1}{2}$.
Multiplication par $(x + 1)$ puis évaluation en $x = -1$ : $\beta = \dfrac{1}{-1 - 1} = -\dfrac{1}{2}$.
Donc
\[ \frac{1}{x^2 - 1} = \frac{1}{2} \left( \frac{1}{x - 1} - \frac{1}{x + 1} \right), \]
et
\[ \int_2^3 \frac{dx}{x^2 - 1} = \frac{1}{2} \left[ \ln|x - 1| - \ln|x + 1| \right]_2^3 = \frac{1}{2} \left( \ln 2 - \ln 4 - \ln 1 + \ln 3 \right) = \frac{1}{2} \ln \frac{3}{2}. \]
La décomposition en éléments simples a transformé un calcul difficile en deux primitives évidentes. Cette section systématise la technique.

\subsection{Théorème de décomposition sur $\mathbb{C}$}

\begin{theoreme}[DES sur $\mathbb{C}(X)$]
Soit $F = A/B \in \mathbb{C}(X)$ sous forme irréductible, avec $B = b\, \prod_{i=1}^{r} (X - a_i)^{m_i}$ (les $a_i$ pôles distincts de multiplicités $m_i$, $b \in \mathbb{C}^*$). Alors $F$ s'écrit de manière \emph{unique} sous la forme
\[ F = E(X) + \sum_{i=1}^{r} \sum_{k=1}^{m_i} \frac{\lambda_{i, k}}{(X - a_i)^k}, \qquad \lambda_{i, k} \in \mathbb{C}, \]
où $E \in \mathbb{C}[X]$ est la partie entière de $F$ (le polynôme tel que $\deg(F - E) < 0$).
\end{theoreme}

\begin{proof}[Idée de la démonstration]
\emph{Existence.} Par récurrence sur le nombre de pôles distincts. On commence par soustraire la partie entière $E$. Pour le pôle $a_1$ de multiplicité $m_1$, on extrait par multiplication par $(X - a_1)^{m_1}$ et division euclidienne tous les termes $\lambda_{1, k}/(X - a_1)^k$ ; la fraction restante a un pôle de moins.

\emph{Unicité.} Si deux décompositions coïncident, leur différence est un polynôme strictement propre nul. L'unicité des coefficients de la « partie singulière » en chaque pôle découle de la minimalité du dénominateur.
\end{proof}

\subsection{Théorème de décomposition sur $\mathbb{R}$}

\begin{theoreme}[DES sur $\mathbb{R}(X)$]
Soit $F = A/B \in \mathbb{R}(X)$ sous forme irréductible. Soit
\[ B = b\, \prod_{i=1}^{r} (X - a_i)^{m_i} \cdot \prod_{j=1}^{s} (X^2 + p_j X + q_j)^{n_j} \]
la factorisation de $B$ en irréductibles réels (où les facteurs de degré $2$ ont un discriminant $p_j^2 - 4 q_j < 0$). Alors $F$ se décompose de manière \emph{unique} :
\[ \boxed{\ F = E(X) + \sum_{i=1}^{r} \sum_{k=1}^{m_i} \frac{\alpha_{i, k}}{(X - a_i)^k} + \sum_{j=1}^{s} \sum_{l=1}^{n_j} \frac{\beta_{j, l}\, X + \gamma_{j, l}}{(X^2 + p_j X + q_j)^l}\ } \]
avec $E \in \mathbb{R}[X]$, $\alpha_{i, k}, \beta_{j, l}, \gamma_{j, l} \in \mathbb{R}$.

Les fractions du premier type sont appelées \emph{éléments simples de première espèce}, celles du second type \emph{éléments simples de seconde espèce}.
\end{theoreme}

\begin{proof}[Idée]
On part de la DES sur $\mathbb{C}(X)$. Les pôles complexes non réels apparaissent par couples conjugués $(a, \bar{a})$ de même multiplicité. On combine les termes $\dfrac{\lambda}{(X - a)^k}$ et $\dfrac{\bar{\lambda}}{(X - \bar{a})^k}$ : leur somme est un polynôme à coefficients réels divisé par $((X - a)(X - \bar{a}))^k = (X^2 - 2 \mathrm{Re}(a) X + |a|^2)^k$. En réduisant dans $\mathbb{R}[X]$, on obtient les éléments simples de seconde espèce.
\end{proof}

\subsection{Méthodes pratiques de calcul des coefficients}

\paragraph{Méthode 1 : multiplication par $(X - a)^m$ pour le coefficient de plus haute multiplicité.}
Soit $a$ pôle de multiplicité $m$ dans la DES, contribuant le terme $\dfrac{\lambda_m}{(X - a)^m} + \cdots + \dfrac{\lambda_1}{X - a}$.
On calcule
\[ \lambda_m = \big[(X - a)^m \cdot F(X)\big]_{X = a}. \]
C'est l'unique coefficient extractible directement par évaluation : les autres exigent un travail supplémentaire (méthode 2 ou 3).

\paragraph{Méthode 2 : valeurs particulières.}
On écrit la DES avec coefficients inconnus, puis l'on évalue $F$ en des valeurs commodes (par exemple $X = 0$ ou $X = 2$ si $0$ et $2$ ne sont pas pôles), pour obtenir un système linéaire.

\paragraph{Méthode 3 : identification par développement asymptotique.}
On multiplie l'identité par le dénominateur, on développe les deux membres comme polynômes, et l'on identifie les coefficients de chaque puissance de $X$.

\paragraph{Méthode 4 : formule du résidu pour pôles simples.}
Si $a$ est pôle simple ($m = 1$) et $F = A/B$ avec $A(a), B'(a) \neq 0$, alors le coefficient de $1/(X - a)$ est :
\[ \lambda_a = \frac{A(a)}{B'(a)}. \]
\textit{Démonstration :} on écrit $B(X) = (X - a) Q(X)$ avec $Q(a) \neq 0$. Alors $B'(a) = Q(a)$ et $\lambda_a = A(a)/Q(a) = A(a)/B'(a)$.

\paragraph{Méthode 5 : dérivation pour les pôles multiples.}
Si $a$ est pôle de multiplicité $m$, en notant $G(X) = (X - a)^m F(X)$, on a
\[ \lambda_{m - k} = \frac{1}{k!}\, G^{(k)}(a) \quad \text{pour } k = 0, 1, \dots, m - 1. \]
C'est la formule de Taylor appliquée à $G$ en $a$.

\subsection{Exemples résolus (Difficulté ascendante)}

\textbf{Exemple 1 (Pôles simples --- Facile).}
Décomposer $F = \dfrac{1}{(X - 1)(X - 2)}$ en éléments simples sur $\mathbb{R}$.

\textit{Correction :} forme cherchée : $F = \dfrac{\alpha}{X - 1} + \dfrac{\beta}{X - 2}$.
Multiplication par $(X - 1)$ et évaluation en $X = 1$ : $\alpha = \dfrac{1}{1 - 2} = -1$.
Multiplication par $(X - 2)$ et évaluation en $X = 2$ : $\beta = \dfrac{1}{2 - 1} = 1$.
\[ F = \frac{-1}{X - 1} + \frac{1}{X - 2}. \]

\medskip

\textbf{Exemple 2 (Pôle multiple --- Intermédiaire).}
Décomposer $F = \dfrac{X + 1}{(X - 1)^2 (X - 2)}$ sur $\mathbb{R}$.

\textit{Correction :} forme cherchée : $F = \dfrac{\alpha_2}{(X - 1)^2} + \dfrac{\alpha_1}{X - 1} + \dfrac{\beta}{X - 2}$.

\smallskip

\emph{Calcul de $\alpha_2$ :} multiplier par $(X - 1)^2$ et évaluer en $X = 1$ :
\[ \alpha_2 = \left[\frac{X + 1}{X - 2}\right]_{X = 1} = \frac{2}{-1} = -2. \]

\emph{Calcul de $\beta$ :} multiplier par $(X - 2)$ et évaluer en $X = 2$ :
\[ \beta = \left[\frac{X + 1}{(X - 1)^2}\right]_{X = 2} = \frac{3}{1} = 3. \]

\emph{Calcul de $\alpha_1$ :} on évalue en une valeur simple, par exemple $X = 0$ :
\[ F(0) = \frac{1}{(-1)^2 (-2)} = -\frac{1}{2}. \]
D'autre part, la DES en $X = 0$ donne :
\[ \frac{\alpha_2}{1} + \frac{\alpha_1}{-1} + \frac{\beta}{-2} = -2 - \alpha_1 - \frac{3}{2} = -\frac{7}{2} - \alpha_1. \]
Identification : $-\dfrac{7}{2} - \alpha_1 = -\dfrac{1}{2}$, soit $\alpha_1 = -3$.

\smallskip

Conclusion :
\[ \frac{X + 1}{(X - 1)^2 (X - 2)} = \frac{-2}{(X - 1)^2} + \frac{-3}{X - 1} + \frac{3}{X - 2}. \]

\medskip

\textbf{Exemple 3 (Élément simple de seconde espèce --- Difficile).}
Décomposer $F = \dfrac{1}{(X^2 + 1)(X - 1)}$ sur $\mathbb{R}$.

\textit{Correction :} le facteur $X^2 + 1$ est irréductible sur $\mathbb{R}$ (discriminant $-4$). La forme cherchée est :
\[ F = \frac{\beta X + \gamma}{X^2 + 1} + \frac{\alpha}{X - 1}. \]

\emph{Calcul de $\alpha$} : multiplication par $(X - 1)$ et évaluation en $X = 1$ :
\[ \alpha = \frac{1}{1 + 1} = \frac{1}{2}. \]

\emph{Calcul de $\beta, \gamma$} : par identification après mise au même dénominateur $(X^2 + 1)(X - 1)$ :
\[ 1 = (\beta X + \gamma)(X - 1) + \alpha (X^2 + 1) = \beta X^2 - \beta X + \gamma X - \gamma + \frac{1}{2} X^2 + \frac{1}{2}. \]
On regroupe :
\[ 1 = \left( \beta + \tfrac{1}{2} \right) X^2 + (\gamma - \beta) X + \left( \tfrac{1}{2} - \gamma \right). \]
Identification :
$\beta + \tfrac{1}{2} = 0$ donne $\beta = -\tfrac{1}{2}$ ;
$\tfrac{1}{2} - \gamma = 1$ donne $\gamma = -\tfrac{1}{2}$ ;
vérification : $\gamma - \beta = -\tfrac{1}{2} - (-\tfrac{1}{2}) = 0$. \checkmark

\smallskip

Conclusion :
\[ \frac{1}{(X^2 + 1)(X - 1)} = \frac{-\tfrac{1}{2} X - \tfrac{1}{2}}{X^2 + 1} + \frac{\tfrac{1}{2}}{X - 1} = \frac{1}{2} \left( \frac{1}{X - 1} - \frac{X + 1}{X^2 + 1} \right). \]

\subsection{Application : calcul de primitives}

Une fois la DES obtenue, l'intégration est immédiate sur chaque terme.

\paragraph{Éléments de première espèce.}
\[ \int \frac{dx}{(x - a)} = \ln |x - a| + C, \qquad \int \frac{dx}{(x - a)^k} = \frac{-1}{(k - 1) (x - a)^{k - 1}} + C \quad (k \ge 2). \]

\paragraph{Éléments de seconde espèce ($x^2 + p x + q$ irréductible).}
On complète le carré : $x^2 + p x + q = (x + p/2)^2 + (q - p^2/4)$, avec $q - p^2/4 > 0$. On pose $u = x + p/2$, et l'on est ramené à $\int \dfrac{\alpha u + \beta}{(u^2 + c^2)^k}\, du$ qui se traite par décomposition $\alpha u du / (u^2 + c^2)^k$ (substitution $v = u^2 + c^2$) et $\beta du / (u^2 + c^2)^k$ (formule de récurrence ou arctan).

\subsection{Exercices pratiques avec Python}
\texttt{sympy.apart} effectue la décomposition en éléments simples automatiquement, sur $\mathbb{Q}$, $\mathbb{R}$ ou $\mathbb{C}$ selon l'extension fournie.

\begin{lstlisting}
import sympy as sp

X = sp.symbols('X')

# 1. DES sur R (default = factorisation sur Q)
F = (X + 1) / ((X - 1)**2 * (X - 2))
print(f"F = {F}")
print(f"  DES sur R : {sp.apart(F, X)}")

# 2. DES avec un facteur de seconde espece
F = 1 / ((X**2 + 1) * (X - 1))
print(f"\nF = {F}")
print(f"  DES sur R : {sp.apart(F, X)}")

# 3. DES sur C : meme fraction, avec extension par I
print(f"  DES sur C : {sp.apart(F, X, full=True).doit()}")

# 4. Application au calcul de primitive
F = 1 / (X**2 - 1)
des = sp.apart(F, X)
print(f"\nF = {F}")
print(f"  DES         : {des}")
print(f"  primitive   : {sp.integrate(F, X)}")

# 5. Verification : la somme des residus est-elle = 0 quand deg F < -1 ?
F = (X**2 + 1) / ((X - 1)*(X - 2)*(X - 3))
des = sp.apart(F, X)
print(f"\nF = {F}, deg F = -1")
print(f"  DES : {des}")
\end{lstlisting}

\subsection{Exercices à faire}
\begin{enumerate}
    \item \textbf{Exercice 1.} Décomposer en éléments simples sur $\mathbb{R}$ :
    \[ F = \frac{1}{(X - 1)(X - 2)(X - 3)}. \]
    \item \textbf{Exercice 2.} Décomposer en éléments simples sur $\mathbb{R}$ :
    \[ F = \frac{X^2 + 2}{(X - 1)^3}. \]
    (Indication : utiliser la dérivation, ou poser $Y = X - 1$ et développer le numérateur en puissances de $Y$.)
    \item \textbf{Exercice 3.} Décomposer en éléments simples sur $\mathbb{R}$ :
    \[ F = \frac{X}{(X^2 + 1)^2}. \]
    \item \textbf{Exercice 4.} En utilisant la décomposition en éléments simples, calculer la primitive
    \[ \int \frac{x}{(x - 1)(x + 2)}\, dx. \]
    \item \textbf{Exercice 5 (synthèse, séries télescopiques).} En décomposant $\dfrac{1}{n(n + 1)}$ en éléments simples, calculer
    \[ S_N = \sum_{n=1}^{N} \frac{1}{n(n + 1)}. \]
    En déduire $\lim_{N \to \infty} S_N$.
\end{enumerate}

\section*{Travaux Pratiques (TP) Python}
\addcontentsline{toc}{section}{Travaux Pratiques (TP) Python}

\vspace{0.3cm}
\noindent L'objectif de ce TP est de construire une bibliothèque \emph{maison} pour les fractions rationnelles : classe \texttt{FractionRationnelle}, mise sous forme irréductible, recherche de pôles et zéros, décomposition en éléments simples, et applications au calcul de primitives et de séries télescopiques. On supposera disponible la classe \texttt{Polynome} construite au TP du chapitre V (ou \texttt{sympy} pour aller plus vite).

\vspace{0.3cm}

\noindent\textbf{Exercice 1 : Classe \texttt{FractionRationnelle}}
\begin{enumerate}
    \item \textbf{[Bloom : Compréhension | Difficulté : Moyenne]} \\
    Créer une classe \texttt{FractionRationnelle} avec deux attributs (\texttt{num}, \texttt{den}) qui sont des objets \texttt{Polynome} (ou des polynômes \texttt{sympy}). Le constructeur doit :
    \begin{itemize}
        \item refuser un dénominateur nul ;
        \item mettre la fraction sous \emph{forme irréductible normalisée} (division par le PGCD, dénominateur unitaire).
    \end{itemize}
    \item \textbf{[Bloom : Application | Difficulté : Moyenne]} \\
    Surcharger \texttt{\_\_repr\_\_} pour un affichage lisible du type \texttt{(num) / (den)}.
\end{enumerate}

\vspace{0.3cm}

\noindent\textbf{Exercice 2 : Arithmétique des fractions rationnelles}
\begin{enumerate}
    \item \textbf{[Bloom : Application | Difficulté : Moyenne]} \\
    Surcharger \texttt{\_\_add\_\_}, \texttt{\_\_sub\_\_}, \texttt{\_\_mul\_\_}, \texttt{\_\_truediv\_\_} pour implémenter $\dfrac{A}{B} + \dfrac{C}{D} = \dfrac{AD + BC}{BD}$ et les autres opérations. Toujours retourner le résultat sous forme irréductible normalisée.
    \item \textbf{[Bloom : Évaluation | Difficulté : Moyenne]} \\
    Vérifier numériquement les axiomes de corps sur 5 fractions aléatoires : $F + (-F) = 0$, $F \cdot (1/F) = 1$, distributivité $F(G + H) = FG + FH$.
\end{enumerate}

\vspace{0.3cm}

\noindent\textbf{Exercice 3 : Partie entière et reste strictement propre}
\begin{enumerate}
    \item \textbf{[Bloom : Application | Difficulté : Moyenne]} \\
    Ajouter une méthode \texttt{partie\_entiere(self)} qui retourne le couple $(E, R/B)$ où $E$ est la partie polynomiale et $R/B$ la partie strictement propre, obtenus par division euclidienne $A = BQ + R$.
    \item \textbf{[Bloom : Évaluation | Difficulté : Moyenne]} \\
    Tester sur $F = \dfrac{X^4 + 1}{X^2 - 1}$ : on doit retrouver $E = X^2 + 1$ et $R/B = \dfrac{2}{X^2 - 1}$.
\end{enumerate}

\vspace{0.3cm}

\noindent\textbf{Exercice 4 : Pôles et zéros avec multiplicités}
\begin{enumerate}
    \item \textbf{[Bloom : Application | Difficulté : Moyenne]} \\
    Ajouter \texttt{poles(self)} et \texttt{zeros(self)} qui retournent un dictionnaire $\{a : m\}$ associant à chaque pôle (resp.\ zéro) sa multiplicité. Travailler sur la forme irréductible.
    \item \textbf{[Bloom : Évaluation | Difficulté : Moyenne]} \\
    Tester sur $F = \dfrac{(X - 2)^3 (X + 1)}{(X - 1)^2 (X + 1)^2}$. Après simplification, identifier zéros et pôles.
    \item \textbf{[Bloom : Synthèse | Difficulté : Moyenne]} \\
    Vérifier numériquement le bilan zéros--pôles dans $\mathbb{C}$ :
    $\sum \mathrm{mult}(z\text{éros}) - \sum \mathrm{mult}(\text{pôles}) = \deg F$.
\end{enumerate}

\vspace{0.3cm}

\noindent\textbf{Exercice 5 : Décomposition en éléments simples (pôles simples)}
\begin{enumerate}
    \item \textbf{[Bloom : Synthèse | Difficulté : Moyenne]} \\
    Soit $F = A/B$ sous forme irréductible avec $B$ scindé à pôles simples sur $\mathbb{C}$. Implémenter \texttt{des\_poles\_simples(F)} qui calcule les coefficients $\lambda_a = \dfrac{A(a)}{B'(a)}$ pour chaque pôle simple $a$.
    \item \textbf{[Bloom : Évaluation | Difficulté : Moyenne]} \\
    Tester sur $F = \dfrac{1}{(X - 1)(X - 2)(X - 3)}$. On doit obtenir $\dfrac{1/2}{X-1} + \dfrac{-1}{X-2} + \dfrac{1/2}{X-3}$.
    \item \textbf{[Bloom : Évaluation | Difficulté : Moyenne]} \\
    Comparer avec \texttt{sympy.apart(F, X)}.
\end{enumerate}

\vspace{0.3cm}

\noindent\textbf{Exercice 6 : DES pour pôle multiple par dérivation (Taylor)}
\begin{enumerate}
    \item \textbf{[Bloom : Synthèse | Difficulté : Difficile]} \\
    Soit $a$ un pôle de multiplicité $m$ de $F = A/B$. En posant $G(X) = (X - a)^m F(X)$ (qui est un polynôme), les coefficients de $1/(X-a)^k$ dans la DES s'obtiennent par les dérivées :
    \[ \lambda_{m - k} = \frac{G^{(k)}(a)}{k!}, \qquad k = 0, 1, \dots, m - 1. \]
    Implémenter \texttt{coefficients\_pole\_multiple(F, a, m)} qui retourne la liste $[\lambda_1, \lambda_2, \dots, \lambda_m]$.
    \item \textbf{[Bloom : Évaluation | Difficulté : Difficile]} \\
    Tester sur $F = \dfrac{X^2 + 2}{(X - 1)^3}$. Calculer la DES complète et vérifier en additionnant les éléments simples.
\end{enumerate}

\vspace{0.3cm}

\noindent\textbf{Exercice 7 : Application --- Sommation de séries télescopiques}
\begin{enumerate}
    \item \textbf{[Bloom : Application | Difficulté : Moyenne]} \\
    Calculer $\displaystyle S_N = \sum_{n=1}^{N} \dfrac{1}{n(n+1)(n+2)}$ par DES. La fraction $\dfrac{1}{X(X+1)(X+2)}$ se décompose en :
    \[ \frac{1}{X(X+1)(X+2)} = \frac{1/2}{X} - \frac{1}{X+1} + \frac{1/2}{X+2}. \]
    Vérifier cette décomposition (par DES Python ou par identification).
    \item \textbf{[Bloom : Synthèse | Difficulté : Difficile]} \\
    En utilisant la décomposition, montrer que la somme se télescope et calculer une formule fermée pour $S_N$. Vérifier sur $N = 10, 100, 1000$ que la formule fermée et la sommation directe coïncident.
    \item \textbf{[Bloom : Évaluation | Difficulté : Moyenne]} \\
    En déduire $\sum_{n=1}^{\infty} \dfrac{1}{n(n+1)(n+2)} = \dfrac{1}{4}$. Vérifier numériquement avec $N = 10^6$.
\end{enumerate}

\vspace{0.3cm}

\noindent\textbf{Exercice 8 : Application --- Calcul de primitives}
\begin{enumerate}
    \item \textbf{[Bloom : Application | Difficulté : Moyenne]} \\
    En utilisant la DES, calculer à la main les primitives suivantes :
    \[ I_1 = \int \frac{dx}{x^2 - 1}, \qquad I_2 = \int \frac{x\, dx}{(x - 1)(x - 2)}, \qquad I_3 = \int \frac{dx}{x(x^2 + 1)}. \]
    \item \textbf{[Bloom : Évaluation | Difficulté : Moyenne]} \\
    Comparer avec \texttt{sympy.integrate} : les expressions doivent coïncider à une constante près (utiliser \texttt{sympy.simplify} sur la différence).
    \item \textbf{[Bloom : Création | Difficulté : Difficile]} \\
    Écrire \texttt{integrer\_fraction\_rationnelle(F, X)} qui automatise le calcul : DES, intégration de chaque élément simple ($\ln$ pour 1$^{\text{re}}$ espèce simple, $\arctan$ + $\ln$ pour 2$^{\text{nde}}$ espèce). Tester sur $\dfrac{1}{(X^2 + 1)(X - 1)}$ et comparer à \texttt{sympy.integrate}.
\end{enumerate}


\end{document}
